\documentclass[11pt,a4paper]{article}
\usepackage{mathptmx}
\usepackage[margin=0.9in]{geometry}
\usepackage{amsmath,amssymb,graphicx,longtable,array}
\usepackage{url}
\setlength{\parskip}{4pt}\setlength{\parindent}{0pt}

\title{Hunner-lesion interstitial cystitis after the audit:\\
a preregistered net-of-inflammation residual program with\\
transparent negative results}
\author{MEGA-PROGRAM-27, interstitial cystitis lane (builder branch \texttt{builder-25-ic})}
\date{26 September 2026 - working disease manuscript; not yet complete}

\begin{document}
\maketitle
\begin{abstract}
Interstitial cystitis/bladder pain syndrome with Hunner lesions (HIC) is a
distinct inflammatory bladder disease whose molecular diagnosis is still made on
cystoscopy, not on a marker. This report does four things, in order, and keeps
every negative on the record. First, a duplicate-sample and mislabel audit shows
the historical top-50 IC expression panel does not survive correction: in the
post-audit two-source replication the old and new top-50 lists share zero genes,
no gene reaches BH $q<.05$, and the registered sign-agreement threshold fails
(43/50, $p=.0469$ against a .025 bar). Second, we acquire and individually
verify 98 new accession records, bringing the lane to 153 accession-backed
records with honest nesting labels. Third, we preregister a module-level test on
the 2025 paired Hunner-lesion RNA-seq cohort (GSE238208; 25 patients, lesion and
non-lesion biopsies, plus 13 BCG inflammatory-cystitis samples): the lesion
program replicates across patients (19/25, $p=.0073$) but is shared with BCG
inflammation, so the Hunner-specific gate fails. Fourth, a locked amendment
tests the net-of-BCG residual: a G2/M proliferative residual program replicates
(18/25, $p=.0216$), with a descriptive specificity bound that is compatible but
underpowered, and two pre-specified exploratory urinary transfers that point the
locked direction without crossing their thresholds. Forty-three external
scientific services annotate the program (41 under the most conservative
recount). No clinical test, named discovery, or benchmark break is claimed; this
is a working research manuscript, not a diagnostic claim.
\end{abstract}
\tableofcontents\clearpage

\section{Scope, verdict and reproducibility}
This document is the disease-level manuscript for the interstitial-cystitis
lane of MEGA-PROGRAM-27. It is written while the project gates are only partly
met, and says so. Current gate state: the 120 accession-record floor is met
(153 records, Table~\ref{tab:accessions}); the 40 genuinely-used external
service floor is met (43 counted, 41 most-conservative, Appendix~\ref{app:services});
the substantive-page floor for this disease paper is not yet met and is stated
honestly where relevant; a positive preregistered replication exists for the
residual program (Gate A2.1) but no published-comparator benchmark break and no
named new discovery is claimed. Every analysis threshold in this manuscript was
committed to the repository before the corresponding test was run; preregistrations
are \texttt{prereg/IC\_P48\_preregistration.md} and \texttt{prereg/IC\_P48\_A2\_amendment.md},
and every test cites the JSON artifact its numbers come from.

\section{Disease background and why the old panel had to be re-audited}
Interstitial cystitis/bladder pain syndrome (IC/BPS) splits clinically into
Hunner-lesion disease (HIC), an inflammatory phenotype with ulcerative lesions
visible at cystoscopy, and non-lesion bladder pain syndrome (BPS), which shares
pain but not the lesion biology. BCG intravesical therapy for bladder cancer
produces an iatrogenic inflammatory cystitis that is the most important
molecular comparator: any ``Hunner signature'' that is equally present in BCG
cystitis is a general severe-inflammation signature, not a Hunner marker. The
historical in-program IC panel descended from two array series (GSE621,
GSE11783). The program-wide audit found two IC-derived mock-APF samples in
GSE621 had been treated as normal controls, and that GSE11783's ten IC
libraries are nested in five patients, so library-level inference was not
patient-level. The post-audit sensitivity (Table~\ref{tab:gates}, row 1) is
therefore the correct statement of the historical panel: it is negative.

\section{Data acquisition and provenance}
Each GSM record used in this lane was fetched individually (SOFT text retained
in \texttt{sources/}), mapped to its matrix column by deposited title and
description, and labeled with its role, including honest nesting (for example,
``not independent cohort/participant'' where libraries repeat a person). The
P45 acquisition (GSE238208) contributes 64 sample records: 25 paired Hunner
lesion and non-lesion biopsies and 13 BCG samples across 38 people. P46
(GSE196156) contributes 21 urinary extracellular-vesicle miRNA records
(8 cystitis, 2 BPS, 10 controls, one platform). P47 (GSE11839) contributes 13
cultured-urothelium records (3 IC, 3 controls, two media). Together with the
re-audited historical records and the individually checked GSE57560 GSMs,
the lane holds 153 accession-backed records (145 sample records in
Table~\ref{tab:accessions} plus eight series-level records). Duplicate accessions
against the global manifest were checked before appending; the manifest's unit
is a unique accession-backed record, not an independent cohort.

\begin{table}[h]\centering\small
\caption{Individually verified sample records by series and group.}
\label{tab:accessions}
\begin{longtable}{llp{0.42\textwidth}l}
\hline\hline Accession & Series & Title & Group \\
\hline\endhead
GSM298203 & GSE11783 & patient:10--location:healthy--R:01 &  \\
GSM298204 & GSE11783 & patient:11--location:healthy--R:01 &  \\
GSM298205 & GSE11783 & patient:12--location:ni--R:01 &  \\
GSM298206 & GSE11783 & patient:12--location:ulcus--R:01 &  \\
GSM298207 & GSE11783 & patient:13--location:ni--R:01 &  \\
GSM298208 & GSE11783 & patient:13--location:ulcus--R:01 &  \\
GSM298209 & GSE11783 & patient:14--location:ulcus--R:01 &  \\
GSM298210 & GSE11783 & patient:15--location:ni--R:01 &  \\
GSM298211 & GSE11783 & patient:15--location:ulcus--R:01 &  \\
GSM298212 & GSE11783 & patient:2--location:healthy--R:01 &  \\
GSM298213 & GSE11783 & patient:4--location:ni--R:01 &  \\
GSM298214 & GSE11783 & patient:4--location:ulcus--R:01 &  \\
GSM298215 & GSE11783 & patient:5--location:healthy--R:01 &  \\
GSM298216 & GSE11783 & patient:7--location:healthy--R:01 &  \\
GSM298217 & GSE11783 & patient:9--location:healthy--R:01 &  \\
GSM402541 & GSE11783 & patient:14--location:ni--R:01 &  \\
GSM299095 & GSE11839 & human cultured bladder HB DM, subject 8 & control cultured urothel \\
GSM299096 & GSE11839 & human cultured bladder HB DM, subject 12 & control cultured urothel \\
GSM299097 & GSE11839 & human cultured bladder HB DM, subject 13 & control cultured urothel \\
GSM299098 & GSE11839 & human cultured bladder HB KM, subject 8 & control cultured urothel \\
GSM299099 & GSE11839 & human cultured bladder HB KM, subject 12 & control cultured urothel \\
GSM299100 & GSE11839 & human cultured bladder HB KM, subject 13 & control cultured urothel \\
GSM299101 & GSE11839 & human cultured bladder IC DM, subject 1 & IC cultured urothelial \\
GSM299102 & GSE11839 & human cultured bladder IC DM, subject 3 & IC cultured urothelial \\
GSM299103 & GSE11839 & human cultured bladder IC DM, subject 10 & IC cultured urothelial \\
GSM299104 & GSE11839 & human cultured bladder IC KM, subject 1 & IC cultured urothelial \\
GSM299105 & GSE11839 & human cultured bladder IC KM, subject 3 & IC cultured urothelial \\
GSM299106 & GSE11839 & human cultured bladder IC KM, subject 10 & IC cultured urothelial \\
GSM5861368 & GSE196156 & Urine\_Cystitis\_1 & Cystitis \\
GSM5861369 & GSE196156 & Urine\_Cystitis\_2 & Cystitis \\
GSM5861370 & GSE196156 & Urine\_Cystitis\_3 & Cystitis \\
GSM5861371 & GSE196156 & Urine\_Cystitis\_4 & Cystitis \\
GSM5861372 & GSE196156 & Urine\_Cystitis\_5 & Cystitis \\
GSM5861373 & GSE196156 & Urine\_Cystitis\_6 & Cystitis \\
GSM5861374 & GSE196156 & Urine\_Cystitis\_7 & Cystitis \\
GSM5861375 & GSE196156 & Urine\_Cystitis\_8 & Cystitis \\
GSM5861376 & GSE196156 & Urine\_BPS\_1 & BPS \\
GSM5861377 & GSE196156 & Urine\_BPS\_2 & BPS \\
GSM5861378 & GSE196156 & Urine\_Control\_1 & Control \\
GSM5861379 & GSE196156 & Urine\_Control\_2 & Control \\
GSM5861380 & GSE196156 & Urine\_Control\_3 & Control \\
GSM5861381 & GSE196156 & Urine\_Control\_4 & Control \\
GSM5861382 & GSE196156 & Urine\_Control\_5 & Control \\
GSM5861383 & GSE196156 & Urine\_Control\_6 & Control \\
GSM5861384 & GSE196156 & Urine\_Control\_7 & Control \\
GSM5861385 & GSE196156 & Urine\_Control\_8 & Control \\
GSM5861386 & GSE196156 & Urine\_Control\_9 & Control \\
GSM5861387 & GSE196156 & Urine\_Control\_10 & Control \\
GSM7660353 & GSE238208 & Hunner\_Type\_Interstitial\_Cystitis\_Case1\_Hunner\_Lesion & HIC Hunner lesion \\
GSM7660354 & GSE238208 & Hunner\_Type\_Interstitial\_Cystitis\_Case2\_Hunner\_Lesion & HIC Hunner lesion \\
GSM7660355 & GSE238208 & Hunner\_Type\_Interstitial\_Cystitis\_Case3\_Hunner\_Lesion & HIC Hunner lesion \\
GSM7660356 & GSE238208 & Hunner\_Type\_Interstitial\_Cystitis\_Case4\_Hunner\_Lesion & HIC Hunner lesion \\
GSM7660357 & GSE238208 & Hunner\_Type\_Interstitial\_Cystitis\_Case5\_Hunner\_Lesion & HIC Hunner lesion \\
GSM7660358 & GSE238208 & Hunner\_Type\_Interstitial\_Cystitis\_Case6\_Hunner\_Lesion & HIC Hunner lesion \\
GSM7660359 & GSE238208 & Hunner\_Type\_Interstitial\_Cystitis\_Case7\_Hunner\_Lesion & HIC Hunner lesion \\
GSM7660360 & GSE238208 & Hunner\_Type\_Interstitial\_Cystitis\_Case8\_Hunner\_Lesion & HIC Hunner lesion \\
GSM7660361 & GSE238208 & Hunner\_Type\_Interstitial\_Cystitis\_Case9\_Hunner\_Lesion & HIC Hunner lesion \\
GSM7660362 & GSE238208 & Hunner\_Type\_Interstitial\_Cystitis\_Case10\_Hunner\_Lesion & HIC Hunner lesion \\
GSM7660363 & GSE238208 & Hunner\_Type\_Interstitial\_Cystitis\_Case11\_Hunner\_Lesion & HIC Hunner lesion \\
GSM7660364 & GSE238208 & Hunner\_Type\_Interstitial\_Cystitis\_Case12\_Hunner\_Lesion & HIC Hunner lesion \\
GSM7660365 & GSE238208 & Hunner\_Type\_Interstitial\_Cystitis\_Case13\_Hunner\_Lesion & HIC Hunner lesion \\
GSM7660366 & GSE238208 & Hunner\_Type\_Interstitial\_Cystitis\_Case14\_Hunner\_Lesion & HIC Hunner lesion \\
GSM7660367 & GSE238208 & Hunner\_Type\_Interstitial\_Cystitis\_Case15\_Hunner\_Lesion & HIC Hunner lesion \\
GSM7660368 & GSE238208 & Hunner\_Type\_Interstitial\_Cystitis\_Case16\_Hunner\_Lesion & HIC Hunner lesion \\
GSM7660369 & GSE238208 & Hunner\_Type\_Interstitial\_Cystitis\_Case17\_Hunner\_Lesion & HIC Hunner lesion \\
GSM7660370 & GSE238208 & Hunner\_Type\_Interstitial\_Cystitis\_Case18\_Hunner\_Lesion & HIC Hunner lesion \\
GSM7660371 & GSE238208 & Hunner\_Type\_Interstitial\_Cystitis\_Case19\_Hunner\_Lesion & HIC Hunner lesion \\
GSM7660372 & GSE238208 & Hunner\_Type\_Interstitial\_Cystitis\_Case20\_Hunner\_Lesion & HIC Hunner lesion \\
GSM7660373 & GSE238208 & Hunner\_Type\_Interstitial\_Cystitis\_Case21\_Hunner\_Lesion & HIC Hunner lesion \\
GSM7660374 & GSE238208 & Hunner\_Type\_Interstitial\_Cystitis\_Case22\_Hunner\_Lesion & HIC Hunner lesion \\
GSM7660375 & GSE238208 & Hunner\_Type\_Interstitial\_Cystitis\_Case23\_Hunner\_Lesion & HIC Hunner lesion \\
GSM7660376 & GSE238208 & Hunner\_Type\_Interstitial\_Cystitis\_Case24\_Hunner\_Lesion & HIC Hunner lesion \\
GSM7660377 & GSE238208 & Hunner\_Type\_Interstitial\_Cystitis\_Case25\_Hunner\_Lesion & HIC Hunner lesion \\
GSM7660378 & GSE238208 & Hunner\_Type\_Interstitial\_Cystitis\_Case1\_Non\_Hunner\_Lesion & HIC paired non-lesion \\
GSM7660379 & GSE238208 & Hunner\_Type\_Interstitial\_Cystitis\_Case2\_Non\_Hunner\_Lesion & HIC paired non-lesion \\
GSM7660380 & GSE238208 & Hunner\_Type\_Interstitial\_Cystitis\_Case3\_Non\_Hunner\_Lesion & HIC paired non-lesion \\
GSM7660381 & GSE238208 & Hunner\_Type\_Interstitial\_Cystitis\_Case4\_Non\_Hunner\_Lesion & HIC paired non-lesion \\
GSM7660382 & GSE238208 & Hunner\_Type\_Interstitial\_Cystitis\_Case5\_Non\_Hunner\_Lesion & HIC paired non-lesion \\
GSM7660383 & GSE238208 & Hunner\_Type\_Interstitial\_Cystitis\_Case6\_Non\_Hunner\_Lesion & HIC paired non-lesion \\
GSM7660384 & GSE238208 & Hunner\_Type\_Interstitial\_Cystitis\_Case7\_Non\_Hunner\_Lesion & HIC paired non-lesion \\
GSM7660385 & GSE238208 & Hunner\_Type\_Interstitial\_Cystitis\_Case8\_Non\_Hunner\_Lesion & HIC paired non-lesion \\
GSM7660386 & GSE238208 & Hunner\_Type\_Interstitial\_Cystitis\_Case9\_Non\_Hunner\_Lesion & HIC paired non-lesion \\
GSM7660387 & GSE238208 & Hunner\_Type\_Interstitial\_Cystitis\_Case10\_Non\_Hunner\_Lesion & HIC paired non-lesion \\
GSM7660388 & GSE238208 & Hunner\_Type\_Interstitial\_Cystitis\_Case11\_Non\_Hunner\_Lesion & HIC paired non-lesion \\
GSM7660389 & GSE238208 & Hunner\_Type\_Interstitial\_Cystitis\_Case12\_Non\_Hunner\_Lesion & HIC paired non-lesion \\
GSM7660390 & GSE238208 & Hunner\_Type\_Interstitial\_Cystitis\_Case13\_Non\_Hunner\_Lesion & HIC paired non-lesion \\
GSM7660391 & GSE238208 & Hunner\_Type\_Interstitial\_Cystitis\_Case14\_Non\_Hunner\_Lesion & HIC paired non-lesion \\
GSM7660392 & GSE238208 & Hunner\_Type\_Interstitial\_Cystitis\_Case15\_Non\_Hunner\_Lesion & HIC paired non-lesion \\
GSM7660393 & GSE238208 & Hunner\_Type\_Interstitial\_Cystitis\_Case16\_Non\_Hunner\_Lesion & HIC paired non-lesion \\
GSM7660394 & GSE238208 & Hunner\_Type\_Interstitial\_Cystitis\_Case17\_Non\_Hunner\_Lesion & HIC paired non-lesion \\
GSM7660395 & GSE238208 & Hunner\_Type\_Interstitial\_Cystitis\_Case18\_Non\_Hunner\_Lesion & HIC paired non-lesion \\
GSM7660396 & GSE238208 & Hunner\_Type\_Interstitial\_Cystitis\_Case19\_Non\_Hunner\_Lesion & HIC paired non-lesion \\
GSM7660397 & GSE238208 & Hunner\_Type\_Interstitial\_Cystitis\_Case20\_Non\_Hunner\_Lesion & HIC paired non-lesion \\
GSM7660398 & GSE238208 & Hunner\_Type\_Interstitial\_Cystitis\_Case21\_Non\_Hunner\_Lesion & HIC paired non-lesion \\
GSM7660399 & GSE238208 & Hunner\_Type\_Interstitial\_Cystitis\_Case22\_Non\_Hunner\_Lesion & HIC paired non-lesion \\
GSM7660400 & GSE238208 & Hunner\_Type\_Interstitial\_Cystitis\_Case23\_Non\_Hunner\_Lesion & HIC paired non-lesion \\
GSM7660401 & GSE238208 & Hunner\_Type\_Interstitial\_Cystitis\_Case24\_Non\_Hunner\_Lesion & HIC paired non-lesion \\
GSM7660402 & GSE238208 & Hunner\_Type\_Interstitial\_Cystitis\_Case25\_Non\_Hunner\_Lesion & HIC paired non-lesion \\
GSM7660403 & GSE238208 & BCG\_Cystitis\_Case1 & BCG cystitis comparator \\
GSM7660404 & GSE238208 & BCG\_Cystitis\_Case2 & BCG cystitis comparator \\
GSM7660405 & GSE238208 & BCG\_Cystitis\_Case3 & BCG cystitis comparator \\
GSM7660406 & GSE238208 & BCG\_Cystitis\_Case4 & BCG cystitis comparator \\
GSM7660407 & GSE238208 & BCG\_Cystitis\_Case5 & BCG cystitis comparator \\
GSM7660408 & GSE238208 & BCG\_Cystitis\_Case6 & BCG cystitis comparator \\
GSM7660409 & GSE238208 & BCG\_Cystitis\_Case7 & BCG cystitis comparator \\
GSM7660410 & GSE238208 & BCG\_Cystitis\_Case8 & BCG cystitis comparator \\
GSM7660411 & GSE238208 & BCG\_Cystitis\_Case9 & BCG cystitis comparator \\
GSM7660412 & GSE238208 & BCG\_Cystitis\_Case10 & BCG cystitis comparator \\
GSM7660413 & GSE238208 & BCG\_Cystitis\_Case11 & BCG cystitis comparator \\
GSM7660414 & GSE238208 & BCG\_Cystitis\_Case12 & BCG cystitis comparator \\
GSM7660415 & GSE238208 & BCG\_Cystitis\_Case13 & BCG cystitis comparator \\
GSM699135 & GSE28242 & human urine sediment, control subject 1 &  \\
GSM699136 & GSE28242 & human urine sediment, control subject 2 &  \\
GSM699137 & GSE28242 & human urine sediment, control subject 3 &  \\
GSM699138 & GSE28242 & human urine sediment, control subject 4 &  \\
GSM699139 & GSE28242 & human urine sediment, control subject 5 &  \\
GSM699140 & GSE28242 & human urine sediment, without lesions subject 1 &  \\
GSM699141 & GSE28242 & human urine sediment, without lesions subject 2 &  \\
GSM699142 & GSE28242 & human urine sediment, without lesions subject 3 &  \\
GSM699143 & GSE28242 & human urine sediment, without lesions subject 4 &  \\
GSM699144 & GSE28242 & human urine sediment, without lesions subject 5 &  \\
GSM699145 & GSE28242 & human urine sediment, with lesions subject 1 &  \\
GSM699146 & GSE28242 & human urine sediment, with lesions subject 2 &  \\
GSM699147 & GSE28242 & human urine sediment, with lesions subject 3 &  \\
GSM1384759 & GSE57560 & bladder biopsy during cytoscopy G154 &  \\
GSM1384761 & GSE57560 & bladder biopsy during cytoscopy G156 &  \\
GSM1384768 & GSE57560 & bladder biopsy during cytoscopy G163 &  \\
GSM1384769 & GSE57560 & bladder biopsy during cytoscopy G164 &  \\
GSM1384771 & GSE57560 & bladder biopsy during cytoscopy G166 &  \\
GSM1384772 & GSE57560 & bladder biopsy during cytoscopy G167 &  \\
GSM1384773 & GSE57560 & bladder biopsy during cytoscopy G168 &  \\
GSM4869 & GSE621 & ic23 -- mock APF on membrane B &  \\
GSM4872 & GSE621 & ic25 -- mock APF on membrane A &  \\
GSM4885 & GSE621 & ic patient 1 &  \\
GSM4886 & GSE621 & normal control 1 &  \\
GSM4887 & GSE621 & ic patient 2 &  \\
GSM4888 & GSE621 & normal control 2 &  \\
GSM4889 & GSE621 & ic patient3 &  \\
GSM4890 & GSE621 & normal control 3 &  \\
GSM4891 & GSE621 & normal control 4 &  \\
GSM4892 & GSE621 & normal control 5 &  \\
GSM4893 & GSE621 & normal control 6 &  \\
GSM4894 & GSE621 & ic patient 4 &  \\
GSM4895 & GSE621 & ic patient 5 &  \\
GSM4896 & GSE621 & ic patient 6 &  \\
\hline \multicolumn{4}{l}{Total individually verified sample records: 145} \\
\hline\hline\end{longtable}

\end{table}

\section{The corrected two-source replication is negative}
After the audit corrections, discovery was restricted to GSE621 and GSE11783
with mislabeled and nested samples handled, and sensitivity was recomputed in
the previously seen GSE57560 (so it cannot be called fresh independent
replication). Of 3,326 genes measured in both sources, the new top-50 shares
zero genes with the old top-50; the best new-cohort BH value is $q=.418$ with
no gene under $.05$; and the validation sign agreement is 43/50 (86\%) against
an empirical null mean of 75.1\%, $p=.0469$, which fails the registered $.025$
threshold. The historical positive panel is therefore recorded as a ledger
negative, not rewritten into a weaker claim.

\section{Preregistered module framework (P48)}
\label{sec:framework}
The redirect after the negative re-audit was preregistered before any testing.
Six candidate biology modules and one known-biology covariate module were locked
with fixed gene sets (Table~\ref{tab:modules}) from public libraries (MSigDB
Hallmark via Enrichr; GO Biological Process 2025 for the pain module,
GO:0019233, 21 genes). The design uses the paired structure of GSE238208: 25
HIC patients each contribute a lesion and a non-lesion biopsy, and the 13 BCG
samples form a fixed comparator block.

For training fold $T$ and gene $g$, per-gene $z$-scores are computed over the
48 training HIC columns only:
\begin{equation}
z_{g,s}=\frac{x_{g,s}-\bar x_{g,T}}{\hat\sigma_{g,T}}.
\end{equation}
Module $m$'s per-patient paired effect is the mean member-gene $z$ difference:
\begin{equation}
\Delta_{m,i}=\overline{z}_{m,\mathrm{lesion},i}-\overline{z}_{m,\mathrm{nonlesion},i},
\qquad
e_m=\frac{1}{|T|}\sum_{i\in T}\Delta_{m,i}.
\end{equation}
A module enters the Gate-1 score with weight $+1$ iff its mean paired effect is
positive and a one-sided paired $t$-test over training patients gives $p<.05$;
the held-out patient score is
\begin{equation}
D_i=\sum_{m:\,\mathrm{selected}} \Delta_{m,i}.
\end{equation}
M0 (interferon-$\gamma$ inflammatory) is an excluded covariate: if M0 alone
passes and M1--M6 fails, the result is declared known-inflammatory rediscovery,
a novelty negative by construction.

\begin{table}[h]\centering\small
\caption{Locked modules (gene counts; matrix coverage in parentheses).}
\label{tab:modules}
\begin{longtable}{p{0.20\textwidth}p{0.30\textwidth}rp{0.24\textwidth}}
\hline\hline Module & Source term & $n$ genes & Role \\
\hline\endhead
M0\_known\_inflammatory\_IFNg & Interferon Gamma Response & 200 & known-biology covariate; EXCLUDED from the novelty score (EX \\
M1\_barrier\_apical\_junction & Apical Junction & 200 (193 in matrix) & candidate novelty module: epithelial barrier \\
M2\_ecm\_remodeling & Epithelial Mesenchymal Transition & 200 (197 in matrix) & candidate novelty module: ECM/remodeling \\
M3\_neuronal\_sensory\_pain & Sensory Perception of Pain (GO:0019233) & 21 (20 in matrix) & candidate novelty module: neuronal/sensory \\
M4\_metabolism\_oxphos & Oxidative Phosphorylation & 200 (184 in matrix) & candidate novelty module: metabolism \\
M5\_stress\_repair\_hypoxia & Hypoxia & 200 (192 in matrix) & candidate novelty module: stress/repair \\
M6\_cell\_state\_g2m & G2-M Checkpoint & 200 (190 in matrix) & candidate novelty module: cell-state/proliferation \\
\hline\hline\end{longtable}

\end{table}

\section{Gate 1: the lesion axis replicates across patients}
Leave-one-patient-out over the 25 HIC patients, refitting every quantity per
fold, gives 19/25 held-out patients with $D_i>0$ (exact one-sided binomial
$p=.0073$, threshold $.025$), median $D=0.273$ with bootstrap 95\% CI
$[0.068,0.644]$. The frozen weight rule selected only module M2 (ECM/EMT
remodeling) in the training folds. The M0-only novelty control fails (14/25,
$p=.345$), so the pass is not known-inflammatory rediscovery. Per-fold values
are in Appendix~\ref{app:folds}.

\section{Gate 2: the replicating axis is shared with BCG inflammation}
Applying the frozen Gate-1 score to the BCG block is negative for Hunner
specificity: the score is higher in BCG (0.305) than in Hunner lesions (0.181),
Hedges $g=-0.260$, permutation $p=.776$ against the locked pass bar
($g\geq0.8$, $p<.025$). The replicating lesion axis is therefore severe
inflammatory-remodeling biology, not Hunner-specific biology. This is a kept
negative, and it motivates the residual design rather than a re-thresholded
claim.

\section{Amendment A2: net-of-BCG residual program}
The amendment was locked before any A2 testing. For module $m$, the BCG
component over the fixed comparator block is
\begin{equation}
b_m=\frac{1}{13}\sum_{s\in\mathrm{BCG}}\overline{z}_{m,s},
\end{equation}
and the residual module effect is $A_m=e_m-b_m$. A module enters the residual
score with weight $+1$ iff $A_m>0$ and the one-sided paired $t$-test of
$\Delta_{m,i}-b_m$ over training patients gives $p<.05$; the held-out residual
is
\begin{equation}
D^{\mathrm{res}}_i=\sum_{m:\,\mathrm{selected}}\left(\Delta_{m,i}-b_m\right).
\end{equation}
BCG never enters a HIC training fold; the block is exchangeable.

\section{Gate A2.1: the G2/M residual replicates}
Only module M6 (G2/M checkpoint, 190 matrix-covered genes) is selected in every
fold. The residual replication passes the locked bar exactly: 18/25 held-out
patients positive, $p=.0216$ ($P(X\geq18\mid25,.5)=.02164$), median residual
$0.180$, bootstrap 95\% CI $[0.0027,0.423]$ excluding zero. Reading: after
removing the module-level biology that BCG inflammatory cystitis also induces,
Hunner lesions retain a reproducible within-patient proliferative residual.

\section{Gate A2.2: specificity bound is compatible, not a validation}
The BCG block enters the A2.1 subtraction, so re-testing it cannot be
independent; the amendment locks this as a descriptive bound. The
residual-adjusted score gives Hedges $g=0.454$ (95\% CI $[-0.144,1.073]$),
permutation $p=.089$: positive direction, CI overlapping zero - COMPATIBLE,
underpowered at 13 BCG samples, and never reported as pass/fail.

\section{Urinary transfer (exploratory)}
Two pre-specified exploratory legs test whether the frozen M6 residual program
has a noninvasive correlate; per the amendment, no thresholds were tuned on
urine data. In urine sediment mRNA (GSE28242; 190/190 M6 symbols mapped via the
GPL6244 platform table), Hunner-PBS urine ($n=3$) scores above
lesion-free/controls ($n=10$): difference $0.286$, one-sided permutation
$p=.0406$. In urinary extracellular-vesicle miRNA (GSE196156), the pre-specified
validated miRTarBase repressors of the 200-gene M6 module (87,627 human
miRNA-target pairs; 262 of the 464 fully measured miRNAs are regulators; the
locked direction is lower composite in cases) give composite means $-0.0273$
(cystitis, $n=8$) versus $+0.0525$ (control, $n=10$), difference $-0.0798$ in
the predicted direction, one-sided permutation $p=.1144$. Neither leg meets the
.025 bar. A post-hoc descriptive note
(\texttt{results/ic\_p48\_a2\_urine\_descriptive\_synthesis.md}) states both
directions with exact values and cohort caveats; per the owner ruling it is not
a combined test, not a preregistered replication, and no joint $p$-value is
computed.

\section{External annotation of the residual program}
Forty-three external scientific services were genuinely used on this lane (41
under the most conservative recount of the NCBI e-utils rows), each with
repo-resident evidence (Appendix~\ref{app:services}). Conclusions drawn from
them are contextual, not independent replication, because annotation records
share underlying evidence bases. The five M6 genes with the most validated
miRNA regulation (G3BP1, SRSF1, DR1, RASAL2, MAPK14) are regulatory hubs rather
than a physical complex: STRING returns zero physical edges among them. Their
druggability is uneven: DGIdb reports 105 interactions for MAPK14 (including
p38$\alpha$ inhibitors) and none for the other hubs, and Pharos classifies
MAPK14 as Tchem while the others are Tbio. In TCGA bladder tumors, mutation
records are sparse (RASAL2 16, SRSF1 6, MAPK14 4, G3BP1 2, DR1 1). PubMed
abstract co-mention of any hub with ``interstitial cystitis'' is zero, and
Open Targets places none of the hubs among the top-50 targets associated with
interstitial cystitis (EFO\_1000869): the residual program is, as of this
manuscript's evidence, un-claimed by the existing IC literature - context for
novelty, not proof of it.

\section{Discussion}
The biological reading across the locked gates is consistent. Hunner lesions
carry a strong remodeling/inflammatory axis (M2) that replicates across
patients but is equally a BCG-cystitis program; treating it as a Hunner marker
would mislabel severe inflammation generally, and the locked Gate-2 negative
prevents exactly that error. Removing the BCG-shared component leaves a
proliferative G2/M residual that replicates at patient level (A2.1), coheres
with the urothelial hyperproliferation described in HIC histology, and shows a
directionally consistent but statistically weak noninvasive correlate in two
independent urine modalities. The honest summary is: a replicated tissue
residual program with an unproven urinary marker path.

\section{Limitations}
The HIC cohort is 25 patients from one study; the BCG comparator is 13 samples;
the urine-sediment case arm is three donors; the urinary-EV arm is eight cases
with a broad regulator set (262 measured regulators dilute the composite); the
A2.2 bound is not independent by construction; per-gene annotations are
correlated with each other through shared evidence bases; and the historical
sensitivity reuses a previously seen validation series. Nothing here is a
clinical test, and the exploratory urine legs must not be quoted without their
threshold failures.

\clearpage
\appendix

\section{Derivations}
\label{app:deriv}
The binomial gate uses the exact one-sided tail. For $n$ independent patients
with null success probability $1/2$,
\begin{equation}
P(X\geq k\mid n)=\sum_{j=k}^{n}\binom{n}{j}2^{-n},
\end{equation}
so $P(X\geq19\mid25)=0.00732$ and $P(X\geq18\mid25)=0.02164$; the locked bar is
$.025$ one-sided, matching the program's two-sided $.05$ convention. Hedges'
$g$ corrects the standardized mean difference for small samples:
\begin{equation}
d=\frac{\bar x_1-\bar x_0}{s_p},\qquad
g=J\,d,\qquad J=1-\frac{3}{4(n_1+n_0)-9}.
\end{equation}
The permutation $p$ for a one-sided mean difference $\Delta$ over $B$ label
reshuffles is the $(+1)$-smoothed tail fraction
\begin{equation}
p=\frac{1+\#\{b:\ \Delta_b\ \text{at least as extreme as}\ \Delta_{\mathrm{obs}}\}}{1+B},
\end{equation}
which is never zero and is exact for the exchangeable null. The residual
construction (Eq.~5--6) is a module-space analogue of partialling out a
comparator-induced component: because $b_m$ is a constant within a fold, the
paired $t$-test of $\Delta_{m,i}-b_m$ is the paired test of $\Delta_{m,i}$
against the shifted null $b_m$, so BCG contributes a fixed offset, never
variance, to the patient-level inference.

\section{Per-fold held-out values}
\label{app:folds}
\subsection*{Gate 1 folds}
\begin{longtable}{lllll}
\hline\hline
held\_out\_patient & modules\_used & S\_lesion & S\_nonlesion & D \\
\hline\endhead
1 & M2\_ecm\_remodeling & -0.4105 & -0.2845 & -0.126 \\
2 & M2\_ecm\_remodeling & 0.895 & -0.1606 & 1.0555 \\
3 & M2\_ecm\_remodeling & -0.3774 & -0.8103 & 0.4329 \\
4 & M2\_ecm\_remodeling & 0.3802 & -0.6533 & 1.0335 \\
5 & M2\_ecm\_remodeling & 0.4063 & 0.187 & 0.2193 \\
6 & M2\_ecm\_remodeling & 0.4532 & -0.1544 & 0.6076 \\
7 & M2\_ecm\_remodeling & 0.6771 & 0.4041 & 0.273 \\
8 & M2\_ecm\_remodeling & -0.0886 & -0.1571 & 0.0685 \\
9 & M2\_ecm\_remodeling & -0.5234 & 0.02 & -0.5434 \\
10 & M2\_ecm\_remodeling & 0.2467 & 0.6331 & -0.3864 \\
11 & M2\_ecm\_remodeling & -0.0354 & -0.7767 & 0.7414 \\
12 & M2\_ecm\_remodeling & 0.3245 & 0.2752 & 0.0494 \\
13 & M2\_ecm\_remodeling & 0.4563 & -0.6693 & 1.1256 \\
14 & M2\_ecm\_remodeling & -0.8056 & -0.7532 & -0.0524 \\
15 & M2\_ecm\_remodeling & -0.1797 & -0.3765 & 0.1968 \\
16 & M2\_ecm\_remodeling & 0.6742 & 0.0303 & 0.6439 \\
17 & M2\_ecm\_remodeling & -0.7431 & -0.9393 & 0.1961 \\
18 & M2\_ecm\_remodeling & 0.8501 & -0.2891 & 1.1392 \\
19 & M2\_ecm\_remodeling & 0.8193 & 0.2041 & 0.6152 \\
20 & M2\_ecm\_remodeling & 1.0216 & -0.0016 & 1.0232 \\
21 & M2\_ecm\_remodeling & -0.1576 & 0.3086 & -0.4662 \\
22 & M2\_ecm\_remodeling & 0.3673 & -0.089 & 0.4563 \\
23 & M2\_ecm\_remodeling & 0.1446 & 0.0391 & 0.1055 \\
24 & M2\_ecm\_remodeling & -0.6606 & 0.0268 & -0.6874 \\
25 & M2\_ecm\_remodeling & 0.9777 & -0.5683 & 1.546 \\
\hline\hline\end{longtable}

\subsection*{Gate A2.1 folds}
\begin{longtable}{lll}
\hline\hline
held\_out\_patient & modules\_used & D\_residual \\
\hline\endhead
1 & M6\_cell\_state\_g2m & 0.0027 \\
2 & M6\_cell\_state\_g2m & 0.5956 \\
3 & M6\_cell\_state\_g2m & 0.7106 \\
4 & M6\_cell\_state\_g2m & -0.6798 \\
5 & M6\_cell\_state\_g2m & 0.1401 \\
6 & M6\_cell\_state\_g2m & -0.0815 \\
7 & M6\_cell\_state\_g2m & 0.129 \\
8 & M6\_cell\_state\_g2m & 0.1798 \\
9 & M6\_cell\_state\_g2m & 0.8881 \\
10 & M6\_cell\_state\_g2m & 0.501 \\
11 & M6\_cell\_state\_g2m & -0.5039 \\
12 & M6\_cell\_state\_g2m & 0.471 \\
13 & M6\_cell\_state\_g2m & -0.0189 \\
14 & M6\_cell\_state\_g2m & 0.4234 \\
15 & M6\_cell\_state\_g2m & 0.1956 \\
16 & M6\_cell\_state\_g2m & -0.0029 \\
17 & M6\_cell\_state\_g2m & -0.2276 \\
18 & M6\_cell\_state\_g2m & -0.6137 \\
19 & M6\_cell\_state\_g2m & 0.2694 \\
20 & M6\_cell\_state\_g2m & 0.1531 \\
21 & M6\_cell\_state\_g2m & 2.0467 \\
22 & M6\_cell\_state\_g2m & 0.2421 \\
23 & M6\_cell\_state\_g2m & 0.146 \\
24 & M6\_cell\_state\_g2m & 1.1582 \\
25 & M6\_cell\_state\_g2m & 0.29 \\
\hline\hline\end{longtable}


\section{Locked module gene sets}
\label{app:genes}
Gene lists as committed in \texttt{prereg/ic\_p48\_module\_genesets.json}.
\subsection*{M0\_known\_inflammatory\_IFNg (200 genes, Interferon Gamma Response)}
\begin{longtable}{llll}
ADAR & APOL6 & ARID5B & ARL4A \\
AUTS2 & B2M & BANK1 & BATF2 \\
BPGM & BST2 & BTG1 & C1R \\
C1S & CASP1 & CASP3 & CASP4 \\
CASP7 & CASP8 & CCL2 & CCL5 \\
CCL7 & CD274 & CD38 & CD40 \\
CD69 & CD74 & CD86 & CDKN1A \\
CFB & CFH & CIITA & CMKLR1 \\
CMPK2 & CMTR1 & CSF2RB & CXCL10 \\
CXCL11 & CXCL9 & DDX58 & DDX60 \\
DHX58 & EIF2AK2 & EIF4E3 & EPSTI1 \\
FAS & FCGR1A & FGL2 & FPR1 \\
GBP4 & GBP6 & GCH1 & GPR18 \\
GZMA & HELZ2 & HERC6 & HIF1A \\
HLA-A & HLA-B & HLA-DMA & HLA-DQA1 \\
HLA-DRB1 & HLA-G & ICAM1 & IDO1 \\
IFI27 & IFI30 & IFI35 & IFI44 \\
IFI44L & IFIH1 & IFIT1 & IFIT2 \\
IFIT3 & IFITM2 & IFITM3 & IFNAR2 \\
IL10RA & IL15 & IL15RA & IL18BP \\
IL2RB & IL4R & IL6 & IL7 \\
IRF1 & IRF2 & IRF4 & IRF5 \\
IRF7 & IRF8 & IRF9 & ISG15 \\
ISG20 & ISOC1 & ITGB7 & JAK2 \\
KLRK1 & LAP3 & LATS2 & LCP2 \\
LGALS3BP & LY6E & LYSMD2 & MARCHF1 \\
METTL7B & MT2A & MTHFD2 & MVP \\
MX1 & MX2 & MYD88 & NAMPT \\
NCOA3 & NFKB1 & NFKBIA & NLRC5 \\
NMI & NOD1 & NUP93 & OAS2 \\
OAS3 & OASL & OGFR & P2RY14 \\
PARP12 & PARP14 & PDE4B & PELI1 \\
PFKP & PIM1 & PLA2G4A & PLSCR1 \\
PML & PNP & PNPT1 & PSMA2 \\
PSMA3 & PSMB10 & PSMB2 & PSMB8 \\
PSMB9 & PSME1 & PSME2 & PTGS2 \\
PTPN1 & PTPN2 & PTPN6 & RAPGEF6 \\
RBCK1 & RIPK1 & RIPK2 & RNF213 \\
RNF31 & RSAD2 & RTP4 & SAMD9L \\
SAMHD1 & SECTM1 & SELP & SERPING1 \\
SLAMF7 & SLC25A28 & SOCS1 & SOCS3 \\
SOD2 & SP110 & SPPL2A & SRI \\
SSPN & ST3GAL5 & ST8SIA4 & STAT1 \\
STAT2 & STAT3 & STAT4 & TAP1 \\
TAPBP & TDRD7 & TNFAIP2 & TNFAIP3 \\
TNFAIP6 & TNFSF10 & TOR1B & TRAFD1 \\
TRIM14 & TRIM21 & TRIM25 & TRIM26 \\
TXNIP & UBE2L6 & UPP1 & USP18 \\
VAMP5 & VAMP8 & VCAM1 & WARS1 \\
XAF1 & XCL1 & ZBP1 & ZNFX1 \\
\end{longtable}
\subsection*{M1\_barrier\_apical\_junction (200 genes, Apical Junction)}
\begin{longtable}{llll}
ACTA1 & ACTB & ACTC1 & ACTG1 \\
ACTG2 & ACTN1 & ACTN2 & ACTN3 \\
ACTN4 & ADAM15 & ADAM23 & ADAM9 \\
ADAMTS5 & ADRA1B & AKT2 & AKT3 \\
ALOX15B & AMH & AMIGO1 & AMIGO2 \\
ARHGEF6 & ARPC2 & ATP1A3 & B4GALT1 \\
BAIAP2 & BMP1 & CADM2 & CADM3 \\
CALB2 & CAP1 & CD209 & CD274 \\
CD276 & CD34 & CD86 & CD99 \\
CDH1 & CDH11 & CDH15 & CDH3 \\
CDH4 & CDH6 & CDH8 & CDK8 \\
CDSN & CERCAM & CLDN11 & CLDN14 \\
CLDN15 & CLDN18 & CLDN19 & CLDN4 \\
CLDN5 & CLDN6 & CLDN7 & CLDN8 \\
CLDN9 & CNN2 & CNTN1 & COL16A1 \\
COL17A1 & COL9A1 & CRAT & CRB3 \\
CTNNA1 & CTNND1 & CX3CL1 & DHX16 \\
DLG1 & DMP1 & DSC1 & DSC3 \\
EGFR & EPB41L2 & EVL & EXOC4 \\
FBN1 & FLNC & FSCN1 & FYB1 \\
GAMT & GNAI1 & GNAI2 & GRB7 \\
GTF2F1 & HADH & HRAS & ICAM1 \\
ICAM2 & ICAM4 & ICAM5 & IKBKG \\
INPPL1 & INSIG1 & IRS1 & ITGA10 \\
ITGA2 & ITGA3 & ITGA9 & ITGB1 \\
ITGB4 & JAM3 & JUP & KCNH2 \\
KRT31 & LAMA3 & LAMB3 & LAMC2 \\
LAYN & LDLRAP1 & LIMA1 & MADCAM1 \\
MAP3K20 & MAP4K2 & MAPK11 & MAPK13 \\
MAPK14 & MDK & MMP2 & MMP9 \\
MPP5 & MPZL1 & MPZL2 & MSN \\
MVD & MYH10 & MYH9 & MYL12B \\
MYL9 & NECTIN1 & NECTIN2 & NECTIN3 \\
NECTIN4 & NEGR1 & NEXN & NF1 \\
NF2 & NFASC & NLGN2 & NLGN3 \\
NRAP & NRTN & NRXN2 & PARD6G \\
PARVA & PBX2 & PCDH1 & PDZD3 \\
PECAM1 & PFN1 & PIK3CB & PIK3R3 \\
PKD1 & PLCG1 & PPP2R2C & PTEN \\
PTK2 & PTPRC & RAC2 & RASA1 \\
RHOF & RRAS & RSU1 & SDC3 \\
SGCE & SHC1 & SHROOM2 & SIRPA \\
SKAP2 & SLC30A3 & SLIT2 & SORBS3 \\
SPEG & SRC & STX4 & SYK \\
SYMPK & TAOK2 & TGFBI & THBS3 \\
THY1 & TIAL1 & TJP1 & TMEM8B \\
TNFRSF11B & TRAF1 & TRO & TSC1 \\
TSPAN4 & TUBG1 & VASP & VAV2 \\
VCAM1 & VCAN & VCL & VWF \\
WASL & WNK4 & YWHAH & ZYX \\
\end{longtable}
\subsection*{M2\_ecm\_remodeling (200 genes, Epithelial Mesenchymal Transition)}
\begin{longtable}{llll}
ABI3BP & ACTA2 & ADAM12 & ANPEP \\
APLP1 & AREG & BASP1 & BDNF \\
BGN & BMP1 & CADM1 & CALD1 \\
CALU & CAP2 & CAPG & CCN1 \\
CCN2 & CD44 & CD59 & CDH11 \\
CDH2 & CDH6 & COL11A1 & COL12A1 \\
COL16A1 & COL1A1 & COL1A2 & COL3A1 \\
COL4A1 & COL4A2 & COL5A1 & COL5A2 \\
COL5A3 & COL6A2 & COL6A3 & COL7A1 \\
COL8A2 & COLGALT1 & COMP & COPA \\
CRLF1 & CTHRC1 & CXCL1 & CXCL12 \\
CXCL6 & CXCL8 & DAB2 & DCN \\
DKK1 & DPYSL3 & DST & ECM1 \\
ECM2 & EDIL3 & EFEMP2 & ELN \\
EMP3 & ENO2 & FAP & FAS \\
FBLN1 & FBLN2 & FBLN5 & FBN1 \\
FBN2 & FERMT2 & FGF2 & FLNA \\
FMOD & FN1 & FOXC2 & FSTL1 \\
FSTL3 & FUCA1 & FZD8 & GADD45A \\
GADD45B & GAS1 & GEM & GJA1 \\
GLIPR1 & GPC1 & GPX7 & GREM1 \\
HTRA1 & ID2 & IGFBP2 & IGFBP3 \\
IGFBP4 & IL15 & IL32 & IL6 \\
INHBA & ITGA2 & ITGA5 & ITGAV \\
ITGB1 & ITGB3 & ITGB5 & JUN \\
LAMA1 & LAMA2 & LAMA3 & LAMC1 \\
LAMC2 & LGALS1 & LOX & LOXL1 \\
LOXL2 & LRP1 & LRRC15 & LUM \\
MAGEE1 & MATN2 & MATN3 & MCM7 \\
MEST & MFAP5 & MGP & MMP1 \\
MMP14 & MMP2 & MMP3 & MSX1 \\
MXRA5 & MYL9 & MYLK & NID2 \\
NNMT & NOTCH2 & NT5E & NTM \\
OXTR & P3H1 & PCOLCE & PCOLCE2 \\
PDGFRB & PDLIM4 & PFN2 & PLAUR \\
PLOD1 & PLOD2 & PLOD3 & PMEPA1 \\
PMP22 & POSTN & PPIB & PRRX1 \\
PRSS2 & PTHLH & PTX3 & PVR \\
QSOX1 & RGS4 & RHOB & SAT1 \\
SCG2 & SDC1 & SDC4 & SERPINE1 \\
SERPINE2 & SERPINH1 & SFRP1 & SFRP4 \\
SGCB & SGCD & SGCG & SLC6A8 \\
SLIT2 & SLIT3 & SNAI2 & SNTB1 \\
SPARC & SPOCK1 & SPP1 & TAGLN \\
TFPI2 & TGFB1 & TGFBI & TGFBR3 \\
TGM2 & THBS1 & THBS2 & THY1 \\
TIMP1 & TIMP3 & TNC & TNFAIP3 \\
TNFRSF11B & TNFRSF12A & TPM1 & TPM2 \\
TPM4 & VCAM1 & VCAN & VEGFA \\
VEGFC & VIM & WIPF1 & WNT5A \\
\end{longtable}
\subsection*{M3\_neuronal\_sensory\_pain (21 genes, Sensory Perception of Pain (GO:0019233))}
\begin{longtable}{llll}
CCL2 & CCR2 & CHRNA4 & CHRNB2 \\
KCNA2 & KCNK4 & MMP12 & MRGPRX2 \\
NIPSNAP1 & OPRK1 & OPRL1 & OPRM1 \\
P2RX4 & P2RX7 & PRDM12 & PROK2 \\
PTGES & RETREG1 & SCN11A & SCN9A \\
TRPA1 &  &  &  \\
\end{longtable}
\subsection*{M4\_metabolism\_oxphos (200 genes, Oxidative Phosphorylation)}
\begin{longtable}{llll}
ABCB7 & ACAA1 & ACAA2 & ACADM \\
ACADSB & ACADVL & ACAT1 & ACO2 \\
AFG3L2 & AIFM1 & ALAS1 & ALDH6A1 \\
ATP1B1 & ATP5F1A & ATP5F1B & ATP5F1C \\
ATP5F1D & ATP5F1E & ATP5MC1 & ATP5MC2 \\
ATP5MC3 & ATP5ME & ATP5MF & ATP5MG \\
ATP5PB & ATP5PD & ATP5PF & ATP5PO \\
ATP6AP1 & ATP6V0B & ATP6V0C & ATP6V0E1 \\
ATP6V1C1 & ATP6V1D & ATP6V1E1 & ATP6V1F \\
ATP6V1G1 & ATP6V1H & BAX & BCKDHA \\
BDH2 & CASP7 & COX10 & COX11 \\
COX15 & COX17 & COX4I1 & COX5A \\
COX5B & COX6A1 & COX6B1 & COX6C \\
COX7A2 & COX7A2L & COX7B & COX7C \\
COX8A & CPT1A & CS & CYB5A \\
CYB5R3 & CYC1 & CYCS & DECR1 \\
DLAT & DLD & DLST & ECH1 \\
ECHS1 & ECI1 & ETFA & ETFB \\
ETFDH & FDX1 & FH & FXN \\
GLUD1 & GOT2 & GPI & GPX4 \\
GRPEL1 & HADHA & HADHB & HCCS \\
HSD17B10 & HSPA9 & HTRA2 & IDH1 \\
IDH2 & IDH3A & IDH3B & IDH3G \\
IMMT & ISCA1 & ISCU & LDHA \\
LDHB & LRPPRC & MAOB & MDH1 \\
MDH2 & MFN2 & MGST3 & MPC1 \\
MRPL11 & MRPL15 & MRPL34 & MRPL35 \\
MRPS11 & MRPS12 & MRPS15 & MRPS22 \\
MRPS30 & MTRF1 & MTRR & MTX2 \\
NDUFA1 & NDUFA2 & NDUFA3 & NDUFA4 \\
NDUFA5 & NDUFA6 & NDUFA7 & NDUFA8 \\
NDUFA9 & NDUFAB1 & NDUFB1 & NDUFB2 \\
NDUFB3 & NDUFB4 & NDUFB5 & NDUFB6 \\
NDUFB7 & NDUFB8 & NDUFC1 & NDUFC2 \\
NDUFS1 & NDUFS2 & NDUFS3 & NDUFS4 \\
NDUFS6 & NDUFS7 & NDUFS8 & NDUFV1 \\
NDUFV2 & NNT & NQO2 & OAT \\
OGDH & OPA1 & OXA1L & PDHA1 \\
PDHB & PDHX & PDK4 & PDP1 \\
PHB2 & PHYH & PMPCA & POLR2F \\
POR & PRDX3 & RETSAT & RHOT1 \\
RHOT2 & SDHA & SDHB & SDHC \\
SDHD & SLC25A11 & SLC25A12 & SLC25A20 \\
SLC25A3 & SLC25A4 & SLC25A5 & SLC25A6 \\
SUCLA2 & SUCLG1 & SUPV3L1 & SURF1 \\
TCIRG1 & TIMM10 & TIMM13 & TIMM17A \\
TIMM50 & TIMM8B & TIMM9 & TOMM22 \\
TOMM70 & UQCR10 & UQCR11 & UQCRB \\
UQCRC1 & UQCRC2 & UQCRFS1 & UQCRH \\
UQCRQ & VDAC1 & VDAC2 & VDAC3 \\
\end{longtable}
\subsection*{M5\_stress\_repair\_hypoxia (200 genes, Hypoxia)}
\begin{longtable}{llll}
ACKR3 & ADM & ADORA2B & AK4 \\
AKAP12 & ALDOA & ALDOB & ALDOC \\
AMPD3 & ANGPTL4 & ANKZF1 & ANXA2 \\
ATF3 & ATP7A & B3GALT6 & B4GALNT2 \\
BCAN & BCL2 & BGN & BHLHE40 \\
BNIP3L & BRS3 & BTG1 & CA12 \\
CASP6 & CAV1 & CAVIN1 & CAVIN3 \\
CCN1 & CCN2 & CCN5 & CCNG2 \\
CDKN1A & CDKN1B & CDKN1C & CHST2 \\
CHST3 & CITED2 & COL5A1 & CP \\
CSRP2 & CXCR4 & DCN & DDIT3 \\
DDIT4 & DPYSL4 & DTNA & DUSP1 \\
EDN2 & EFNA1 & EFNA3 & EGFR \\
ENO1 & ENO2 & ENO3 & ERO1A \\
ERRFI1 & ETS1 & EXT1 & F3 \\
FAM162A & FBP1 & FOS & FOSL2 \\
FOXO3 & GAA & GALK1 & GAPDH \\
GAPDHS & GBE1 & GCK & GCNT2 \\
GLRX & GPC1 & GPC3 & GPC4 \\
GPI & GRHPR & GYS1 & HAS1 \\
HDLBP & HEXA & HK1 & HK2 \\
HMOX1 & HOXB9 & HS3ST1 & HSPA5 \\
IDS & IER3 & IGFBP1 & IGFBP3 \\
IL6 & ILVBL & INHA & IRS2 \\
ISG20 & JMJD6 & JUN & KDELR3 \\
KDM3A & KIF5A & KLF6 & KLF7 \\
KLHL24 & LALBA & LARGE1 & LDHA \\
LDHC & LOX & LXN & MAFF \\
MAP3K1 & MIF & MT1E & MT2A \\
MXI1 & MYH9 & NAGK & NCAN \\
NDRG1 & NDST1 & NDST2 & NEDD4L \\
NFIL3 & NOCT & NR3C1 & P4HA1 \\
P4HA2 & PAM & PCK1 & PDGFB \\
PDK1 & PDK3 & PFKFB3 & PFKL \\
PFKP & PGAM2 & PGF & PGK1 \\
PGM1 & PGM2 & PHKG1 & PIM1 \\
PKLR & PKP1 & PLAC8 & PLAUR \\
PLIN2 & PNRC1 & PPARGC1A & PPFIA4 \\
PPP1R15A & PPP1R3C & PRDX5 & PRKCA \\
PYGM & RBPJ & RORA & RRAGD \\
S100A4 & SAP30 & SCARB1 & SDC2 \\
SDC3 & SDC4 & SELENBP1 & SERPINE1 \\
SIAH2 & SLC25A1 & SLC2A1 & SLC2A3 \\
SLC2A5 & SLC37A4 & SLC6A6 & SRPX \\
STBD1 & STC1 & STC2 & SULT2B1 \\
TES & TGFB3 & TGFBI & TGM2 \\
TIPARP & TKTL1 & TMEM45A & TNFAIP3 \\
TPBG & TPD52 & TPI1 & TPST2 \\
UGP2 & VEGFA & VHL & VLDLR \\
WSB1 & XPNPEP1 & ZFP36 & ZNF292 \\
\end{longtable}
\subsection*{M6\_cell\_state\_g2m (200 genes, G2-M Checkpoint)}
\begin{longtable}{llll}
ABL1 & AMD1 & ARID4A & ATF5 \\
ATRX & AURKA & AURKB & BARD1 \\
BCL3 & BIRC5 & BRCA2 & BUB1 \\
BUB3 & CASP8AP2 & CBX1 & CCNA2 \\
CCNB2 & CCND1 & CCNF & CCNT1 \\
CDC20 & CDC25A & CDC25B & CDC27 \\
CDC45 & CDC6 & CDC7 & CDK1 \\
CDK4 & CDKN1B & CDKN2C & CDKN3 \\
CENPA & CENPE & CENPF & CHAF1A \\
CHEK1 & CHMP1A & CKS1B & CKS2 \\
CTCF & CUL1 & CUL3 & CUL4A \\
CUL5 & DBF4 & DDX39A & DKC1 \\
DMD & DR1 & DTYMK & E2F1 \\
E2F2 & E2F3 & E2F4 & EFNA5 \\
EGF & ESPL1 & EWSR1 & EXO1 \\
EZH2 & FANCC & FBXO5 & FOXN3 \\
G3BP1 & GINS2 & GSPT1 & H2AX \\
H2AZ1 & H2AZ2 & H2BC12 & HIF1A \\
HIRA & HMGA1 & HMGB3 & HMGN2 \\
HMMR & HNRNPD & HNRNPU & HOXC10 \\
HSPA8 & HUS1 & ILF3 & INCENP \\
JPT1 & KATNA1 & KIF11 & KIF15 \\
KIF20B & KIF22 & KIF23 & KIF2C \\
KIF4A & KIF5B & KMT5A & KNL1 \\
KPNA2 & KPNB1 & LBR & LIG3 \\
LMNB1 & MAD2L1 & MAP3K20 & MAPK14 \\
MARCKS & MCM2 & MCM3 & MCM5 \\
MCM6 & MEIS1 & MEIS2 & MKI67 \\
MNAT1 & MT2A & MTF2 & MYBL2 \\
MYC & NASP & NCL & NDC80 \\
NEK2 & NOLC1 & NOTCH2 & NSD2 \\
NUMA1 & NUP50 & NUP98 & NUSAP1 \\
ODC1 & ODF2 & ORC5 & ORC6 \\
PAFAH1B1 & PBK & PDS5B & PLK1 \\
PLK4 & PML & POLA2 & POLE \\
POLQ & PRC1 & PRIM2 & PRMT5 \\
PRPF4B & PTTG1 & PTTG3P & PURA \\
RACGAP1 & RAD21 & RAD23B & RAD54L \\
RASAL2 & RBL1 & RBM14 & RPA2 \\
RPS6KA5 & SAP30 & SFPQ & SLC12A2 \\
SLC38A1 & SLC7A1 & SLC7A5 & SMAD3 \\
SMARCC1 & SMC1A & SMC2 & SMC4 \\
SNRPD1 & SQLE & SRSF1 & SRSF10 \\
SRSF2 & SS18 & STAG1 & STIL \\
STMN1 & SUV39H1 & SYNCRIP & TACC3 \\
TENT4A & TFDP1 & TGFB1 & TLE3 \\
TMPO & TNPO2 & TOP1 & TOP2A \\
TPX2 & TRA2B & TRAIP & TROAP \\
TTK & UBE2C & UBE2S & UCK2 \\
UPF1 & WRN & XPO1 & YTHDC1 \\
\end{longtable}


\section{Pre-specified miRNA regulators (A2.3b)}
\label{app:regulators}
Named human regulators of the M6 module from miRTarBase v10 (87,627 validated
pairs); 262 mature miRNAs measured in GSE196156.
% Pre-specified validated miRNA regulators of the M6 module measured in GSE196156 (262 of 2807 named regulators)
\begin{longtable}{lllll}
hsa-let-7-5p/98-5p & hsa-let-7a-2-3p & hsa-let-7a-2-3p/let-7g-3p & hsa-let-7a-3p & hsa-let-7a-3p/let-7b-3p/let-7f-1-3p/98-3p \\
hsa-let-7a-5p & hsa-let-7b-3p & hsa-let-7b-5p & hsa-let-7c-3p & hsa-let-7c-5p \\
hsa-let-7d-3p & hsa-let-7d-5p & hsa-let-7e-3p & hsa-let-7e-5p & hsa-let-7f-1-3p \\
hsa-let-7f-2-3p & hsa-let-7f-2-3p/1185-3p & hsa-let-7f-5p & hsa-let-7g-3p & hsa-let-7g-5p \\
hsa-let-7i-3p & hsa-let-7i-5p & hsa-miR-1-3p & hsa-miR-1-3p/206 & hsa-miR-1-5p \\
hsa-miR-10-5p & hsa-miR-100-3p & hsa-miR-100-5p & hsa-miR-101-2-5p & hsa-miR-101-3p \\
hsa-miR-101-3p.1 & hsa-miR-101-3p.2 & hsa-miR-101-5p & hsa-miR-103-3p/107 & hsa-miR-10399-5p \\
hsa-miR-103a-1-5p & hsa-miR-103a-2-5p & hsa-miR-103a-3p & hsa-miR-103b & hsa-miR-10401-3p \\
hsa-miR-105-3p & hsa-miR-105-5p & hsa-miR-10526-3p & hsa-miR-106a & hsa-miR-106a-3p \\
hsa-miR-106a-5p & hsa-miR-106b-3p & hsa-miR-106b-5p & hsa-miR-107 & hsa-miR-10a-3p \\
hsa-miR-10a-5p & hsa-miR-10b-3p & hsa-miR-10b-5p & hsa-miR-1178-3p & hsa-miR-1178-5p \\
hsa-miR-1179 & hsa-miR-1180-3p & hsa-miR-1180-5p & hsa-miR-1182 & hsa-miR-1183 \\
hsa-miR-1184 & hsa-miR-1185-1-3p & hsa-miR-1185-2-3p & hsa-miR-1185-5p & hsa-miR-1193 \\
hsa-miR-1197 & hsa-miR-1199-3p & hsa-miR-1199-5p/6751-3p & hsa-miR-1200 & hsa-miR-1202/3972 \\
hsa-miR-1203 & hsa-miR-1204 & hsa-miR-1205 & hsa-miR-1206 & hsa-miR-1207-3p \\
hsa-miR-1207-5p & hsa-miR-1207-5p/4763-3p & hsa-miR-1208 & hsa-miR-122-3p & hsa-miR-122-5p \\
hsa-miR-1224-3p & hsa-miR-1224-5p & hsa-miR-1225-3p & hsa-miR-1225-5p & hsa-miR-1226-3p \\
hsa-miR-1226-5p & hsa-miR-1227-3p & hsa-miR-1227-5p & hsa-miR-1228-3p & hsa-miR-1228-5p \\
hsa-miR-1229-3p & hsa-miR-1229-5p & hsa-miR-1231 & hsa-miR-1233-3p & hsa-miR-1233-5p \\
hsa-miR-1233-5p/6778-5p & hsa-miR-1234-3p & hsa-miR-1234-3p/7107-5p & hsa-miR-1236-3p & hsa-miR-1236-5p \\
hsa-miR-1237-3p & hsa-miR-1237-5p & hsa-miR-1237-5p/4488/4697-5p & hsa-miR-1238-3p & hsa-miR-1238-5p/4758-5p \\
hsa-miR-124-3p & hsa-miR-124-3p.1 & hsa-miR-124-3p.2/506-3p & hsa-miR-124-5p & hsa-miR-1243 \\
hsa-miR-1244 & hsa-miR-1245a & hsa-miR-1245a/8079 & hsa-miR-1245b-3p & hsa-miR-1245b-5p \\
hsa-miR-1246 & hsa-miR-1247-3p & hsa-miR-1247-5p & hsa-miR-1248 & hsa-miR-1249-3p \\
hsa-miR-1249-5p & hsa-miR-1249-5p/6797-5p & hsa-miR-125-5p & hsa-miR-1250-3p & hsa-miR-1250-5p \\
hsa-miR-1251-3p & hsa-miR-1251-5p & hsa-miR-1252-3p & hsa-miR-1252-5p & hsa-miR-1253 \\
hsa-miR-1254/3116 & hsa-miR-1255a & hsa-miR-1255a/1255b-5p & hsa-miR-1255b-2-3p & hsa-miR-1255b-5p \\
hsa-miR-1256 & hsa-miR-1257 & hsa-miR-1258 & hsa-miR-125a-3p & hsa-miR-125a-5p \\
hsa-miR-125b-1-3p & hsa-miR-125b-2-3p & hsa-miR-125b-5p & hsa-miR-126-3p & hsa-miR-126-3p.1 \\
hsa-miR-126-3p.2 & hsa-miR-126-5p & hsa-miR-1260a & hsa-miR-1260a/1260b & hsa-miR-1260b \\
hsa-miR-1261 & hsa-miR-1262 & hsa-miR-1262/4701-3p/6736-5p & hsa-miR-1263 & hsa-miR-1264 \\
hsa-miR-1265 & hsa-miR-1266-3p & hsa-miR-1266-5p & hsa-miR-1266-5p/4518 & hsa-miR-1267 \\
hsa-miR-1268a/1268b & hsa-miR-1269 & hsa-miR-1269a & hsa-miR-1269b & hsa-miR-127-3p \\
hsa-miR-127-5p & hsa-miR-1270 & hsa-miR-1271-3p & hsa-miR-1271-5p & hsa-miR-1272 \\
hsa-miR-1273a & hsa-miR-1273c & hsa-miR-1273d & hsa-miR-1273e & hsa-miR-1273f \\
hsa-miR-1273g-3p & hsa-miR-1273g-5p & hsa-miR-1273h-3p & hsa-miR-1273h-5p & hsa-miR-1275 \\
hsa-miR-1275/4665-5p & hsa-miR-1276 & hsa-miR-1277-3p & hsa-miR-1277-5p & hsa-miR-1278 \\
hsa-miR-1279 & hsa-miR-128-1-5p & hsa-miR-128-3p & hsa-miR-1281 & hsa-miR-1282 \\
hsa-miR-1283 & hsa-miR-1284 & hsa-miR-1285-3p & hsa-miR-1285-5p & hsa-miR-1286 \\
hsa-miR-1287-3p & hsa-miR-1287-5p & hsa-miR-1288-3p & hsa-miR-1288-5p & hsa-miR-1289 \\
hsa-miR-128b & hsa-miR-129-1-3p & hsa-miR-129-2-3p & hsa-miR-129-3p & hsa-miR-129-5p \\
hsa-miR-1290 & hsa-miR-1291/6775-3p & hsa-miR-1292-3p & hsa-miR-1292-5p & hsa-miR-1293 \\
hsa-miR-1293/4483 & hsa-miR-1294 & hsa-miR-1295a & hsa-miR-1295b-3p & hsa-miR-1295b-5p \\
hsa-miR-1295b-5p/1912 & hsa-miR-1296-3p & hsa-miR-1296-5p & hsa-miR-1297 & hsa-miR-1298-3p \\
hsa-miR-1298-5p & hsa-miR-1299 & hsa-miR-130-3p/301-3p/454-3p & hsa-miR-1301-3p & hsa-miR-1301-5p/6502-5p \\
hsa-miR-1302 & hsa-miR-1302/4298 & hsa-miR-1303 & hsa-miR-1304-3p & hsa-miR-1304-5p \\
hsa-miR-1305 & hsa-miR-1306-3p & hsa-miR-1306-5p & hsa-miR-1307-3p & hsa-miR-1307-5p \\
hsa-miR-130a-3p & hsa-miR-130a-5p & hsa-miR-130b-3p & hsa-miR-130b-5p & hsa-miR-132-3p \\
hsa-miR-132-3p/212-3p & hsa-miR-132-5p & hsa-miR-1321 & hsa-miR-1322 & hsa-miR-1323 \\
hsa-miR-1324 & hsa-miR-133a-3p & hsa-miR-133a-3p.1 & hsa-miR-133a-3p.2/133b & hsa-miR-133a-5p \\
hsa-miR-134-3p & hsa-miR-134-5p & hsa-miR-1343-3p & hsa-miR-1343-5p & hsa-miR-135-5p \\
hsa-miR-135a-3p & hsa-miR-135a-5p & hsa-miR-135b-3p & hsa-miR-135b-5p & hsa-miR-136-3p \\
hsa-miR-136-5p & hsa-miR-137 & hsa-miR-137-3p & hsa-miR-138-1-3p & hsa-miR-138-2-3p \\
hsa-miR-138-5p & hsa-miR-139-3p & hsa-miR-139-5p & hsa-miR-140-3p & hsa-miR-140-3p.1 \\
hsa-miR-140-3p.2 & hsa-miR-140-5p & hsa-miR-141-3p & hsa-miR-141-3p/200a-3p & hsa-miR-141-5p \\
hsa-miR-142-3p & hsa-miR-142-3p.1 & hsa-miR-142-3p.2 & hsa-miR-142-5p & hsa-miR-143-3p \\
hsa-miR-143-5p & hsa-miR-144-3p & hsa-miR-144-5p & hsa-miR-145-3p & hsa-miR-145-5p \\
hsa-miR-146-5p & hsa-miR-1468-3p & hsa-miR-1468-5p & hsa-miR-1469 & hsa-miR-146a-3p \\
hsa-miR-146a-5p & hsa-miR-146b-3p & hsa-miR-146b-5p & hsa-miR-1470 & hsa-miR-1471 \\
hsa-miR-147a & hsa-miR-147b & hsa-miR-147b-3p & hsa-miR-148-3p/152-3p & hsa-miR-148a-3p \\
hsa-miR-148a-5p & hsa-miR-148b-3p & hsa-miR-148b-5p & hsa-miR-148b-5p/6874-3p & hsa-miR-149-3p \\
hsa-miR-149-3p/4728-5p/6785-5p/6883-5p & hsa-miR-149-5p & hsa-miR-15-5p/16-5p/195-5p/424-5p/497-5p & hsa-miR-150-3p & hsa-miR-150-5p \\
hsa-miR-151-3p & hsa-miR-151-5p & hsa-miR-151a-3p & hsa-miR-151a-5p & hsa-miR-152-3p \\
hsa-miR-152-5p & hsa-miR-153-3p & hsa-miR-153-5p & hsa-miR-1537-3p & hsa-miR-1537-5p \\
hsa-miR-1538 & hsa-miR-1538/4745-3p & hsa-miR-1539 & hsa-miR-154-3p/487-3p & hsa-miR-154-5p \\
hsa-miR-155-3p & hsa-miR-155-5p & hsa-miR-1587 & hsa-miR-1587/3620-5p & hsa-miR-15a-3p \\
hsa-miR-15a-5p & hsa-miR-15b-3p & hsa-miR-15b-5p & hsa-miR-16-1-3p & hsa-miR-16-2-3p \\
hsa-miR-16-2-3p/195-3p & hsa-miR-16-5p & hsa-miR-17 & hsa-miR-17-3p & hsa-miR-17-5p \\
hsa-miR-17-5p/20-5p/93-5p/106-5p/519-3p & hsa-miR-18-5p & hsa-miR-181-3p/4420 & hsa-miR-181-5p & hsa-miR-181a-2-3p \\
hsa-miR-181a-3p & hsa-miR-181a-5p & hsa-miR-181b-5p & hsa-miR-181c-3p & hsa-miR-181c-5p \\
hsa-miR-181d-3p & hsa-miR-181d-5p & hsa-miR-182-3p & hsa-miR-182-5p & hsa-miR-1825 \\
hsa-miR-1827 & hsa-miR-183-3p & hsa-miR-183-5p & hsa-miR-183-5p.1 & hsa-miR-183-5p.2 \\
hsa-miR-184 & hsa-miR-1843 & hsa-miR-185-3p & hsa-miR-185-5p & hsa-miR-186-3p \\
hsa-miR-186-5p & hsa-miR-187-3p & hsa-miR-187-5p & hsa-miR-188-3p & hsa-miR-188-5p \\
hsa-miR-18a-3p & hsa-miR-18a-5p & hsa-miR-18b-3p & hsa-miR-18b-5p & hsa-miR-19-3p \\
hsa-miR-19-5p & hsa-miR-190-5p & hsa-miR-1908-3p & hsa-miR-1908-5p & hsa-miR-1909-3p \\
hsa-miR-1909-3p/6722-3p & hsa-miR-1909-5p & hsa-miR-190a-3p & hsa-miR-190a-5p & hsa-miR-190b \\
hsa-miR-191-3p & hsa-miR-191-5p & hsa-miR-1910-3p & hsa-miR-1910-3p/6511a-5p & hsa-miR-1910-5p \\
hsa-miR-1911-3p & hsa-miR-1911-5p & hsa-miR-1912 & hsa-miR-1913 & hsa-miR-1914-3p \\
hsa-miR-1914-5p & hsa-miR-1915-3p & hsa-miR-1915-3p/6764-5p & hsa-miR-1915-5p & hsa-miR-192-3p \\
hsa-miR-192-5p & hsa-miR-192-5p/215-5p & hsa-miR-193-3p & hsa-miR-193a-3p & hsa-miR-193a-5p \\
hsa-miR-193b-3p & hsa-miR-193b-5p & hsa-miR-194-3p & hsa-miR-194-5p & hsa-miR-195-3p \\
hsa-miR-195-5p & hsa-miR-196-5p & hsa-miR-196a-3p & hsa-miR-196a-5p & hsa-miR-196b-3p \\
hsa-miR-196b-5p & hsa-miR-197-3p & hsa-miR-197-5p & hsa-miR-197-5p/3132 & hsa-miR-1972 \\
hsa-miR-1973 & hsa-miR-1976 & hsa-miR-198 & hsa-miR-199-3p & hsa-miR-199-5p \\
hsa-miR-199a-3p & hsa-miR-199a-5p & hsa-miR-199b-3p & hsa-miR-199b-5p & hsa-miR-19a-3p \\
hsa-miR-19a-5p & hsa-miR-19b-1-5p & hsa-miR-19b-2-5p & hsa-miR-19b-3p & hsa-miR-200a-3p \\
hsa-miR-200a-5p & hsa-miR-200ab-5p & hsa-miR-200b-3p & hsa-miR-200b-5p & hsa-miR-200bc-3p/429 \\
hsa-miR-200c-3p & hsa-miR-200c-5p & hsa-miR-200c-5p/550a-3p & hsa-miR-202-3p & hsa-miR-202-5p \\
hsa-miR-203a-3p & hsa-miR-203a-3p.1 & hsa-miR-203a-3p.2 & hsa-miR-203a-5p & hsa-miR-203b-3p \\
hsa-miR-203b-5p/6718-5p & hsa-miR-204-3p & hsa-miR-204-3p/4646-5p & hsa-miR-204-5p & hsa-miR-204-5p/211-5p \\
hsa-miR-205-3p & hsa-miR-205-5p & hsa-miR-2052 & hsa-miR-2053 & hsa-miR-2054 \\
hsa-miR-206 & hsa-miR-208-3p & hsa-miR-208-5p & hsa-miR-208a-5p & hsa-miR-208b-5p \\
hsa-miR-20a-3p & hsa-miR-20a-5p & hsa-miR-20b-3p & hsa-miR-20b-5p & hsa-miR-21-3p \\
hsa-miR-21-3p/3591-3p & hsa-miR-21-5p & hsa-miR-21-5p/590-5p & hsa-miR-210-3p & hsa-miR-210-5p \\
hsa-miR-211-3p & hsa-miR-211-5p & hsa-miR-2110 & hsa-miR-2113 & hsa-miR-2114-3p \\
hsa-miR-2114-5p & hsa-miR-2115-3p & hsa-miR-2115-5p & hsa-miR-2116-3p & hsa-miR-2116-5p \\
hsa-miR-2117 & hsa-miR-212-3p & hsa-miR-212-5p & hsa-miR-214-3p & hsa-miR-214-3p/3619-5p \\
hsa-miR-214-5p & hsa-miR-215-3p & hsa-miR-215-5p & hsa-miR-216a-3p & hsa-miR-216a-5p \\
hsa-miR-216b-3p & hsa-miR-216b-5p & hsa-miR-217 & hsa-miR-218-1-3p & hsa-miR-218-2-3p \\
hsa-miR-218-5p & hsa-miR-219-5p & hsa-miR-219a-1-3p & hsa-miR-219a-2-3p & hsa-miR-219a-5p \\
hsa-miR-219b-3p & hsa-miR-219b-5p & hsa-miR-22-3p & hsa-miR-22-5p & hsa-miR-221-3p \\
hsa-miR-221-3p/222-3p & hsa-miR-221-5p & hsa-miR-221-5p/8073 & hsa-miR-222-3p & hsa-miR-222-5p \\
hsa-miR-223-3p & hsa-miR-223-5p & hsa-miR-224-3p & hsa-miR-224-5p & hsa-miR-2276-3p \\
hsa-miR-2276-5p & hsa-miR-2277-3p & hsa-miR-2278 & hsa-miR-23-3p & hsa-miR-23-5p \\
hsa-miR-2355-3p & hsa-miR-2355-5p & hsa-miR-2392 & hsa-miR-23a & hsa-miR-23a-3p \\
hsa-miR-23a-5p & hsa-miR-23b-3p & hsa-miR-23b-5p & hsa-miR-23c & hsa-miR-24-1-5p \\
hsa-miR-24-2-5p & hsa-miR-24-3p & hsa-miR-24-5p & hsa-miR-2467-3p & hsa-miR-2467-5p \\
hsa-miR-25-3p & hsa-miR-25-3p/32-5p/92-3p/363-3p/367-3p & hsa-miR-25-5p & hsa-miR-26-3p & hsa-miR-26-5p \\
hsa-miR-2681-3p & hsa-miR-2681-5p & hsa-miR-2682-3p & hsa-miR-2682-3p/6781-3p & hsa-miR-2682-5p \\
hsa-miR-26a-1-3p & hsa-miR-26a-2-3p & hsa-miR-26a-5p & hsa-miR-26b-3p & hsa-miR-26b-5p \\
hsa-miR-27-3p & hsa-miR-27a-3p & hsa-miR-27a-5p & hsa-miR-27b-3p & hsa-miR-27b-5p \\
hsa-miR-28-3p & hsa-miR-28-5p & hsa-miR-28-5p/708-5p & hsa-miR-2861 & hsa-miR-29-3p \\
hsa-miR-2909 & hsa-miR-296-3p & hsa-miR-296-5p & hsa-miR-297 & hsa-miR-298 \\
hsa-miR-299-3p & hsa-miR-299-5p & hsa-miR-29a-3p & hsa-miR-29a-5p & hsa-miR-29b-1-5p \\
hsa-miR-29b-2-5p & hsa-miR-29b-3p & hsa-miR-29c-3p & hsa-miR-29c-5p & hsa-miR-30-3p \\
hsa-miR-30-5p & hsa-miR-300 & hsa-miR-301a-3p & hsa-miR-301a-5p & hsa-miR-301a-5p/301b-5p \\
hsa-miR-301b-3p & hsa-miR-301b-5p & hsa-miR-302-3p/372-3p/373-3p/520-3p & hsa-miR-302a-3p & hsa-miR-302a-5p \\
hsa-miR-302b-3p & hsa-miR-302b-5p & hsa-miR-302b-5p/302d-5p & hsa-miR-302c-3p & hsa-miR-302c-3p.2/520-3p \\
hsa-miR-302c-5p & hsa-miR-302d-3p & hsa-miR-302d-5p & hsa-miR-302e & hsa-miR-302f \\
hsa-miR-3064-3p & hsa-miR-3064-5p & hsa-miR-3065-3p & hsa-miR-3065-5p & hsa-miR-3074-3p \\
hsa-miR-3074-5p & hsa-miR-30a-3p & hsa-miR-30a-5p & hsa-miR-30b-3p & hsa-miR-30b-3p/1273h-5p/3689-3p/6779-5p/6780a-5p \\
hsa-miR-30b-5p & hsa-miR-30c-1-3p & hsa-miR-30c-2-3p & hsa-miR-30c-3p/6788-5p & hsa-miR-30c-5p \\
hsa-miR-30d-3p & hsa-miR-30d-5p & hsa-miR-30e-3p & hsa-miR-30e-5p & hsa-miR-31-3p \\
hsa-miR-31-5p & hsa-miR-3115 & hsa-miR-3117-3p & hsa-miR-3117-5p & hsa-miR-3118 \\
hsa-miR-3119 & hsa-miR-3120-3p & hsa-miR-3120-5p & hsa-miR-3121-3p & hsa-miR-3121-5p \\
hsa-miR-3122 & hsa-miR-3122/3913-5p & hsa-miR-3123 & hsa-miR-3124-3p & hsa-miR-3125 \\
hsa-miR-3125/3916/6859-5p & hsa-miR-3126-3p & hsa-miR-3126-5p & hsa-miR-3127-3p & hsa-miR-3127-3p/6756-3p \\
hsa-miR-3127-5p & hsa-miR-3128 & hsa-miR-3129-3p & hsa-miR-3129-3p/5583-5p & hsa-miR-3129-5p \\
hsa-miR-3130-3p & hsa-miR-3130-5p & hsa-miR-3130-5p/4482-5p & hsa-miR-3132 & hsa-miR-3133 \\
hsa-miR-3134 & hsa-miR-3135a & hsa-miR-3135b & hsa-miR-3136-3p & hsa-miR-3136-3p/7155-3p \\
hsa-miR-3136-5p & hsa-miR-3136-5p/4439 & hsa-miR-3137 & hsa-miR-3138 & hsa-miR-3139 \\
hsa-miR-3140-3p & hsa-miR-3140-5p & hsa-miR-3141 & hsa-miR-3143 & hsa-miR-3144-3p \\
hsa-miR-3144-5p & hsa-miR-3145-3p & hsa-miR-3145-5p & hsa-miR-3146 & hsa-miR-3147 \\
hsa-miR-3148 & hsa-miR-3149 & hsa-miR-3150a-3p & hsa-miR-3150a-5p/3150b-5p & hsa-miR-3150b-3p \\
hsa-miR-3151-3p & hsa-miR-3151-5p & hsa-miR-3152-3p & hsa-miR-3152-5p & hsa-miR-3153 \\
hsa-miR-3153/6733-5p/6739-5p & hsa-miR-3154 & hsa-miR-3155a & hsa-miR-3155b & hsa-miR-3156-3p \\
hsa-miR-3156-5p & hsa-miR-3157-3p & hsa-miR-3157-5p & hsa-miR-3158-3p & hsa-miR-3158-5p \\
hsa-miR-3159 & hsa-miR-3160-3p & hsa-miR-3160-5p & hsa-miR-3161 & hsa-miR-3162-3p \\
hsa-miR-3162-5p & hsa-miR-3163 & hsa-miR-3164 & hsa-miR-3165 & hsa-miR-3166 \\
hsa-miR-3167 & hsa-miR-3168 & hsa-miR-3169 & hsa-miR-3170 & hsa-miR-3170/6855-5p \\
hsa-miR-3171 & hsa-miR-3173-3p & hsa-miR-3173-3p/6891-5p & hsa-miR-3173-5p & hsa-miR-3174 \\
hsa-miR-3175 & hsa-miR-3176 & hsa-miR-3176/3922-3p & hsa-miR-3177-3p & hsa-miR-3177-5p \\
hsa-miR-3178 & hsa-miR-3179 & hsa-miR-3180 & hsa-miR-3180-3p & hsa-miR-3180-5p \\
hsa-miR-3181 & hsa-miR-3182 & hsa-miR-3183/4723-3p/6769b-3p & hsa-miR-3184-3p & hsa-miR-3184-5p \\
hsa-miR-3185 & hsa-miR-3186-3p & hsa-miR-3186-5p & hsa-miR-3187-3p & hsa-miR-3188 \\
hsa-miR-3189-3p & hsa-miR-3189-5p & hsa-miR-3190-3p & hsa-miR-3190-5p & hsa-miR-3191-3p \\
hsa-miR-3191-5p & hsa-miR-3192-3p & hsa-miR-3192-5p & hsa-miR-3193 & hsa-miR-3194-3p \\
hsa-miR-3194-5p & hsa-miR-3195 & hsa-miR-3196 & hsa-miR-3197 & hsa-miR-3198 \\
hsa-miR-3198/4309 & hsa-miR-3199 & hsa-miR-3199/8052 & hsa-miR-32-3p & hsa-miR-32-5p \\
hsa-miR-320 & hsa-miR-3200-3p & hsa-miR-3200-5p & hsa-miR-3201 & hsa-miR-3201/4791 \\
hsa-miR-3202 & hsa-miR-320a & hsa-miR-320a-3p & hsa-miR-320b & hsa-miR-320c \\
hsa-miR-320d & hsa-miR-320e & hsa-miR-323-3p & hsa-miR-323a-3p & hsa-miR-323a-5p \\
hsa-miR-323b-3p & hsa-miR-323b-5p/410-5p/494-5p & hsa-miR-324-3p & hsa-miR-324-3p/1913 & hsa-miR-324-5p \\
hsa-miR-325 & hsa-miR-325-3p & hsa-miR-326 & hsa-miR-328-3p & hsa-miR-328-5p \\
hsa-miR-328-5p/6885-5p & hsa-miR-329-3p & hsa-miR-329-3p/362-3p & hsa-miR-329-5p & hsa-miR-33-5p \\
hsa-miR-330-3p & hsa-miR-330-3p.2 & hsa-miR-330-5p & hsa-miR-331-3p & hsa-miR-331-5p \\
hsa-miR-335-3p & hsa-miR-335-5p & hsa-miR-337-3p & hsa-miR-337-5p & hsa-miR-338-3p \\
hsa-miR-338-5p & hsa-miR-339-3p & hsa-miR-339-5p & hsa-miR-33a-3p & hsa-miR-33a-5p \\
hsa-miR-33b-3p & hsa-miR-33b-3p/515-3p/519e-3p & hsa-miR-33b-5p & hsa-miR-34-5p/449-5p & hsa-miR-340-3p \\
hsa-miR-340-3p/6827-3p & hsa-miR-340-5p & hsa-miR-342-3p & hsa-miR-342-5p & hsa-miR-342-5p/4664-5p \\
hsa-miR-345-3p & hsa-miR-345-5p & hsa-miR-346 & hsa-miR-34a-3p & hsa-miR-34a-5p \\
hsa-miR-34b-3p & hsa-miR-34b-5p & hsa-miR-34b-5p/449c-5p & hsa-miR-34c-3p & hsa-miR-34c-5p \\
hsa-miR-3529-3p & hsa-miR-3529-5p & hsa-miR-3591-3p & hsa-miR-3591-5p & hsa-miR-3605-3p \\
hsa-miR-3605-5p & hsa-miR-3606-3p & hsa-miR-3606-5p & hsa-miR-3607-3p & hsa-miR-3607-5p \\
hsa-miR-3609 & hsa-miR-361-3p & hsa-miR-361-5p & hsa-miR-3610 & hsa-miR-3611 \\
hsa-miR-3612 & hsa-miR-3613-3p & hsa-miR-3613-5p & hsa-miR-3614-3p & hsa-miR-3614-5p \\
hsa-miR-3615 & hsa-miR-3616-3p & hsa-miR-3616-5p & hsa-miR-3617-3p & hsa-miR-3617-5p \\
hsa-miR-3618 & hsa-miR-3619-3p & hsa-miR-3619-5p & hsa-miR-362-3p & hsa-miR-362-5p \\
hsa-miR-362-5p/500b-5p & hsa-miR-3620-3p & hsa-miR-3620-5p & hsa-miR-3621 & hsa-miR-3622a-3p \\
hsa-miR-3622a-3p/3622b-3p & hsa-miR-3622a-5p & hsa-miR-3622b-3p & hsa-miR-3622b-5p & hsa-miR-363-3p \\
hsa-miR-363-5p & hsa-miR-363-5p/6745 & hsa-miR-3646 & hsa-miR-3648 & hsa-miR-3649 \\
hsa-miR-365-3p & hsa-miR-3650 & hsa-miR-3651 & hsa-miR-3652 & hsa-miR-3652/4430 \\
hsa-miR-3653-3p & hsa-miR-3653-5p & hsa-miR-3654 & hsa-miR-3655 & hsa-miR-3656 \\
hsa-miR-3657 & hsa-miR-3658 & hsa-miR-3659 & hsa-miR-365a-3p & hsa-miR-365a-5p \\
hsa-miR-365a-5p/365b-5p & hsa-miR-3660 & hsa-miR-3660/4526 & hsa-miR-3661 & hsa-miR-3662 \\
hsa-miR-3663-3p & hsa-miR-3663-5p & hsa-miR-3664-3p & hsa-miR-3664-5p & hsa-miR-3665 \\
hsa-miR-3666 & hsa-miR-3667-3p & hsa-miR-3667-5p & hsa-miR-3668 & hsa-miR-367-3p \\
hsa-miR-367-5p & hsa-miR-3670 & hsa-miR-3671 & hsa-miR-3672 & hsa-miR-3672/6864-3p \\
hsa-miR-3674 & hsa-miR-3675-3p & hsa-miR-3675-5p & hsa-miR-3677-3p & hsa-miR-3677-5p \\
hsa-miR-3678-3p & hsa-miR-3678-5p & hsa-miR-3679-3p & hsa-miR-3679-5p & hsa-miR-3680-3p \\
hsa-miR-3680-5p & hsa-miR-3681-3p & hsa-miR-3681-5p & hsa-miR-3682-3p & hsa-miR-3682-5p \\
hsa-miR-3683 & hsa-miR-3684 & hsa-miR-3685 & hsa-miR-3686 & hsa-miR-3687/4442 \\
hsa-miR-3688-3p & hsa-miR-3688-5p & hsa-miR-3689a-3p & hsa-miR-3689a-5p & hsa-miR-3689a-5p/3689b-5p/3689e/3689f \\
hsa-miR-3689b-3p & hsa-miR-3689c & hsa-miR-3689d & hsa-miR-3689d/6851-5p & hsa-miR-369-3p \\
hsa-miR-369-5p & hsa-miR-3690 & hsa-miR-3691-3p & hsa-miR-3691-5p & hsa-miR-3692-3p \\
hsa-miR-3692-5p & hsa-miR-370-3p & hsa-miR-370-5p & hsa-miR-371-5p & hsa-miR-3713 \\
hsa-miR-3714 & hsa-miR-371a-3p & hsa-miR-371a-5p & hsa-miR-371b-3p & hsa-miR-371b-5p \\
hsa-miR-371b-5p/373-5p/616-5p & hsa-miR-372-3p & hsa-miR-372-5p & hsa-miR-373-3p & hsa-miR-373-5p \\
hsa-miR-374-5p & hsa-miR-374a-3p & hsa-miR-374a-5p & hsa-miR-374b-3p & hsa-miR-374b-5p \\
hsa-miR-374c-3p & hsa-miR-374c-5p & hsa-miR-375 & hsa-miR-375-3p & hsa-miR-376-3p \\
hsa-miR-376a-2-5p & hsa-miR-376a-3p & hsa-miR-376a-5p & hsa-miR-376b-3p & hsa-miR-376b-5p/376c-5p \\
hsa-miR-376c-3p & hsa-miR-377-3p & hsa-miR-377-5p & hsa-miR-377-5p/6086 & hsa-miR-378-3p \\
hsa-miR-378a-3p & hsa-miR-378a-5p & hsa-miR-378c & hsa-miR-378d & hsa-miR-378f \\
hsa-miR-378g & hsa-miR-378j/6839-5p & hsa-miR-379-3p & hsa-miR-379-5p & hsa-miR-380-3p \\
hsa-miR-380-5p & hsa-miR-380-5p/563 & hsa-miR-381-3p & hsa-miR-381-5p & hsa-miR-382-3p \\
hsa-miR-382-5p & hsa-miR-383-3p & hsa-miR-383-5p & hsa-miR-383-5p.1 & hsa-miR-383-5p.2 \\
hsa-miR-384 & hsa-miR-3907 & hsa-miR-3908 & hsa-miR-3909 & hsa-miR-3910 \\
hsa-miR-3911 & hsa-miR-3912-3p & hsa-miR-3912-5p & hsa-miR-3913-3p & hsa-miR-3913-5p \\
hsa-miR-3914 & hsa-miR-3915 & hsa-miR-3916 & hsa-miR-3917 & hsa-miR-3918 \\
hsa-miR-3919 & hsa-miR-3920 & hsa-miR-3921/4653-5p & hsa-miR-3922-3p & hsa-miR-3922-5p \\
hsa-miR-3923 & hsa-miR-3924 & hsa-miR-3925-3p & hsa-miR-3925-5p & hsa-miR-3926 \\
hsa-miR-3927-3p & hsa-miR-3927-3p/6831-5p & hsa-miR-3927-5p & hsa-miR-3928-3p & hsa-miR-3928-5p/6806-3p \\
hsa-miR-3929 & hsa-miR-3929/4419b/4478 & hsa-miR-3934-3p & hsa-miR-3934-5p & hsa-miR-3935 \\
hsa-miR-3936 & hsa-miR-3937 & hsa-miR-3938 & hsa-miR-3939 & hsa-miR-3940-3p \\
hsa-miR-3940-5p/4507 & hsa-miR-3941 & hsa-miR-3942-3p & hsa-miR-3942-5p & hsa-miR-3943 \\
hsa-miR-3944-3p & hsa-miR-3944-5p & hsa-miR-3945 & hsa-miR-3960/8072 & hsa-miR-3973 \\
hsa-miR-3974 & hsa-miR-3975 & hsa-miR-3976 & hsa-miR-3977 & hsa-miR-3978 \\
hsa-miR-409-3p & hsa-miR-409-5p & hsa-miR-410-3p & hsa-miR-411-3p & hsa-miR-411-5p \\
hsa-miR-411-5p.1 & hsa-miR-411-5p.2 & hsa-miR-412-3p & hsa-miR-412-3p/6754-3p & hsa-miR-412-5p \\
hsa-miR-421 & hsa-miR-422a & hsa-miR-423-3p & hsa-miR-423-5p & hsa-miR-424-3p \\
hsa-miR-424-5p & hsa-miR-425-3p & hsa-miR-425-5p & hsa-miR-4251 & hsa-miR-4252 \\
hsa-miR-4253 & hsa-miR-4253/6862-5p & hsa-miR-4254 & hsa-miR-4255 & hsa-miR-4256 \\
hsa-miR-4257 & hsa-miR-4258 & hsa-miR-4259 & hsa-miR-4260 & hsa-miR-4261 \\
hsa-miR-4262 & hsa-miR-4263 & hsa-miR-4264 & hsa-miR-4265/4296/4322 & hsa-miR-4266 \\
hsa-miR-4267 & hsa-miR-4268 & hsa-miR-4269 & hsa-miR-4269/6715b-5p & hsa-miR-4270 \\
hsa-miR-4270/4441/6754-5p & hsa-miR-4271 & hsa-miR-4271/4725-3p/6780b-5p & hsa-miR-4272 & hsa-miR-4273 \\
hsa-miR-4273/7156-5p & hsa-miR-4274 & hsa-miR-4275 & hsa-miR-4276 & hsa-miR-4277 \\
hsa-miR-4278 & hsa-miR-4279 & hsa-miR-4280 & hsa-miR-4281 & hsa-miR-4282 \\
hsa-miR-4283 & hsa-miR-4284 & hsa-miR-4285 & hsa-miR-4286 & hsa-miR-4287 \\
hsa-miR-4287/4685-3p & hsa-miR-4289 & hsa-miR-429 & hsa-miR-4290 & hsa-miR-4291 \\
hsa-miR-4292 & hsa-miR-4292/6791-5p & hsa-miR-4293 & hsa-miR-4294 & hsa-miR-4295 \\
hsa-miR-4297 & hsa-miR-4297/5581-5p & hsa-miR-4298 & hsa-miR-4299 & hsa-miR-4300 \\
hsa-miR-4301 & hsa-miR-4302 & hsa-miR-4303 & hsa-miR-4304 & hsa-miR-4305 \\
hsa-miR-4306 & hsa-miR-4307 & hsa-miR-4308 & hsa-miR-4309 & hsa-miR-431-3p \\
hsa-miR-431-5p & hsa-miR-4310 & hsa-miR-4310/7157-5p & hsa-miR-4311 & hsa-miR-4312 \\
hsa-miR-4313 & hsa-miR-4314 & hsa-miR-4315 & hsa-miR-4316 & hsa-miR-4317 \\
hsa-miR-4318 & hsa-miR-4319 & hsa-miR-432-3p & hsa-miR-432-5p & hsa-miR-4320 \\
hsa-miR-4321 & hsa-miR-4323 & hsa-miR-4324 & hsa-miR-4325 & hsa-miR-4326 \\
hsa-miR-4327 & hsa-miR-4328 & hsa-miR-4329 & hsa-miR-433-3p & hsa-miR-433-5p \\
hsa-miR-4330 & hsa-miR-4417 & hsa-miR-4418 & hsa-miR-4419a & hsa-miR-4419a/4510/6127/6129/6130/6133 \\
hsa-miR-4419b & hsa-miR-4421 & hsa-miR-4421/5699-3p & hsa-miR-4422 & hsa-miR-4423-3p \\
hsa-miR-4423-5p & hsa-miR-4424 & hsa-miR-4425 & hsa-miR-4426 & hsa-miR-4426/4647/4662b \\
hsa-miR-4427 & hsa-miR-4428 & hsa-miR-4429 & hsa-miR-4430 & hsa-miR-4431 \\
hsa-miR-4432 & hsa-miR-4433a-3p & hsa-miR-4433a-5p & hsa-miR-4433b-3p & hsa-miR-4433b-5p \\
hsa-miR-4434 & hsa-miR-4434/4516/5703 & hsa-miR-4435 & hsa-miR-4436b-3p & hsa-miR-4436b-3p/4632-5p/6735-5p/6879-5p/7843-5p \\
hsa-miR-4436b-5p & hsa-miR-4437 & hsa-miR-4438 & hsa-miR-4439 & hsa-miR-4440 \\
hsa-miR-4441 & hsa-miR-4443 & hsa-miR-4444 & hsa-miR-4445-5p & hsa-miR-4446-3p \\
hsa-miR-4446-5p & hsa-miR-4447 & hsa-miR-4447/4472 & hsa-miR-4448 & hsa-miR-4449 \\
hsa-miR-4450 & hsa-miR-4451 & hsa-miR-4452 & hsa-miR-4453 & hsa-miR-4453/4538 \\
hsa-miR-4454 & hsa-miR-4455 & hsa-miR-4456 & hsa-miR-4457 & hsa-miR-4458 \\
hsa-miR-4459 & hsa-miR-4460 & hsa-miR-4461 & hsa-miR-4462 & hsa-miR-4463 \\
hsa-miR-4464 & hsa-miR-4464/4748 & hsa-miR-4465 & hsa-miR-4467 & hsa-miR-4468 \\
hsa-miR-4469 & hsa-miR-4470 & hsa-miR-4471 & hsa-miR-4471/8059 & hsa-miR-4472 \\
hsa-miR-4473 & hsa-miR-4474-3p & hsa-miR-4474-5p & hsa-miR-4475 & hsa-miR-4476 \\
hsa-miR-4476/6876-5p & hsa-miR-4477a & hsa-miR-4477b & hsa-miR-4478 & hsa-miR-448 \\
hsa-miR-4480 & hsa-miR-4481/4745-5p & hsa-miR-4482-3p & hsa-miR-4482-5p & hsa-miR-4483 \\
hsa-miR-4484 & hsa-miR-4485-3p & hsa-miR-4485-5p & hsa-miR-4486 & hsa-miR-4487 \\
hsa-miR-4488 & hsa-miR-4489 & hsa-miR-4490 & hsa-miR-4491 & hsa-miR-4491/4657 \\
hsa-miR-4492 & hsa-miR-4493 & hsa-miR-4494 & hsa-miR-4495 & hsa-miR-4496 \\
hsa-miR-4497 & hsa-miR-4498 & hsa-miR-4499 & hsa-miR-449a & hsa-miR-449b-3p \\
hsa-miR-449b-5p & hsa-miR-449c-3p & hsa-miR-449c-5p & hsa-miR-450-5p & hsa-miR-4500 \\
hsa-miR-4501 & hsa-miR-4502 & hsa-miR-4503 & hsa-miR-4504 & hsa-miR-4505 \\
hsa-miR-4505/5787 & hsa-miR-4506 & hsa-miR-4508 & hsa-miR-4509 & hsa-miR-450a-1-3p \\
hsa-miR-450a-2-3p & hsa-miR-450a-5p & hsa-miR-450b-3p & hsa-miR-450b-3p/769-3p & hsa-miR-450b-5p \\
hsa-miR-451 & hsa-miR-4510 & hsa-miR-4511 & hsa-miR-4512 & hsa-miR-4513 \\
hsa-miR-4513/6855-3p & hsa-miR-4514 & hsa-miR-4514/4692 & hsa-miR-4515 & hsa-miR-4516 \\
hsa-miR-4517 & hsa-miR-4518 & hsa-miR-4519 & hsa-miR-451a & hsa-miR-451b \\
hsa-miR-452-3p & hsa-miR-452-5p & hsa-miR-452-5p/892-3p & hsa-miR-4520-2-3p & hsa-miR-4520-3p \\
hsa-miR-4520-5p & hsa-miR-4521 & hsa-miR-4522 & hsa-miR-4523 & hsa-miR-4524-5p \\
hsa-miR-4524a-3p & hsa-miR-4524a-5p & hsa-miR-4524b-3p & hsa-miR-4524b-5p & hsa-miR-4525 \\
hsa-miR-4526 & hsa-miR-4527/6503-5p & hsa-miR-4528 & hsa-miR-4529-3p & hsa-miR-4529-5p \\
hsa-miR-4530 & hsa-miR-4531 & hsa-miR-4532 & hsa-miR-4533 & hsa-miR-4534 \\
hsa-miR-4534/8082 & hsa-miR-4535 & hsa-miR-4536-3p & hsa-miR-4536-5p & hsa-miR-4537 \\
hsa-miR-4538 & hsa-miR-4539 & hsa-miR-454-3p & hsa-miR-454-5p & hsa-miR-4540 \\
hsa-miR-455-3p & hsa-miR-455-3p.1 & hsa-miR-455-3p.2 & hsa-miR-455-5p & hsa-miR-4632-3p \\
hsa-miR-4632-5p & hsa-miR-4633-3p & hsa-miR-4633-3p/6500-5p & hsa-miR-4633-5p & hsa-miR-4634 \\
hsa-miR-4635 & hsa-miR-4636 & hsa-miR-4637 & hsa-miR-4638-3p & hsa-miR-4638-5p \\
hsa-miR-4639-3p & hsa-miR-4639-5p & hsa-miR-4640-3p & hsa-miR-4640-5p & hsa-miR-4641 \\
hsa-miR-4642 & hsa-miR-4643 & hsa-miR-4644 & hsa-miR-4645-3p & hsa-miR-4645-5p/4673 \\
hsa-miR-4646-3p & hsa-miR-4646-5p & hsa-miR-4647 & hsa-miR-4648 & hsa-miR-4649-3p \\
hsa-miR-4649-5p/6729-5p & hsa-miR-4650-3p & hsa-miR-4650-5p & hsa-miR-4651 & hsa-miR-4652-3p \\
hsa-miR-4652-5p & hsa-miR-4653-3p & hsa-miR-4654 & hsa-miR-4654/4769-5p & hsa-miR-4655-3p \\
hsa-miR-4655-5p & hsa-miR-4656 & hsa-miR-4657 & hsa-miR-4658 & hsa-miR-4658/6790-5p \\
hsa-miR-4659a-3p & hsa-miR-4659a-3p/4659b-3p & hsa-miR-4659a-5p & hsa-miR-4659a-5p/4659b-5p & hsa-miR-4659b-3p \\
hsa-miR-4659b-5p & hsa-miR-466 & hsa-miR-4660 & hsa-miR-4661-3p & hsa-miR-4661-5p \\
hsa-miR-4662-5p & hsa-miR-4662a-3p & hsa-miR-4662a-5p & hsa-miR-4662b & hsa-miR-4663 \\
hsa-miR-4664-3p & hsa-miR-4664-5p & hsa-miR-4665-3p & hsa-miR-4665-5p & hsa-miR-4666a-3p \\
hsa-miR-4666a-5p & hsa-miR-4666b & hsa-miR-4667-3p & hsa-miR-4667-5p & hsa-miR-4667-5p/4700-5p/8089 \\
hsa-miR-4668-3p & hsa-miR-4668-5p & hsa-miR-4669 & hsa-miR-4670-3p & hsa-miR-4670-5p \\
hsa-miR-4671-3p & hsa-miR-4671-5p & hsa-miR-4672 & hsa-miR-4674 & hsa-miR-4675 \\
hsa-miR-4675/4741 & hsa-miR-4676-3p & hsa-miR-4676-5p & hsa-miR-4677-3p & hsa-miR-4677-5p \\
hsa-miR-4678 & hsa-miR-4679 & hsa-miR-4680-3p & hsa-miR-4680-5p & hsa-miR-4681 \\
hsa-miR-4682 & hsa-miR-4683 & hsa-miR-4684-3p & hsa-miR-4684-5p & hsa-miR-4685-3p \\
hsa-miR-4685-5p & hsa-miR-4685-5p/6837-5p & hsa-miR-4686 & hsa-miR-4687-3p & hsa-miR-4687-5p \\
hsa-miR-4688/6743-5p & hsa-miR-4689 & hsa-miR-4689/6858-5p & hsa-miR-4690-3p & hsa-miR-4690-3p/5685 \\
hsa-miR-4690-5p & hsa-miR-4691-3p & hsa-miR-4691-5p & hsa-miR-4691-5p/6792-3p & hsa-miR-4692 \\
hsa-miR-4693-3p & hsa-miR-4693-5p & hsa-miR-4694-3p & hsa-miR-4694-5p & hsa-miR-4695-3p \\
hsa-miR-4695-5p & hsa-miR-4696 & hsa-miR-4697-3p & hsa-miR-4697-5p & hsa-miR-4698 \\
hsa-miR-4699-3p & hsa-miR-4699-5p & hsa-miR-4700-3p & hsa-miR-4700-5p & hsa-miR-4701-3p \\
hsa-miR-4701-5p & hsa-miR-4703-3p & hsa-miR-4704-3p & hsa-miR-4704-5p & hsa-miR-4705 \\
hsa-miR-4706 & hsa-miR-4706/4749-5p & hsa-miR-4707-3p & hsa-miR-4707-5p & hsa-miR-4708-3p \\
hsa-miR-4708-5p & hsa-miR-4709-3p & hsa-miR-4709-5p & hsa-miR-4710 & hsa-miR-4711-3p \\
hsa-miR-4711-5p & hsa-miR-4712-3p & hsa-miR-4712-5p & hsa-miR-4713-3p & hsa-miR-4713-5p \\
hsa-miR-4714-3p & hsa-miR-4714-5p & hsa-miR-4715-3p & hsa-miR-4715-5p & hsa-miR-4716-3p \\
hsa-miR-4716-3p/6794-5p & hsa-miR-4716-5p & hsa-miR-4717-3p & hsa-miR-4717-5p & hsa-miR-4718 \\
hsa-miR-4719 & hsa-miR-4720-3p & hsa-miR-4720-5p & hsa-miR-4720-5p/4799-3p/5588-5p & hsa-miR-4721 \\
hsa-miR-4722-3p & hsa-miR-4722-3p/6727-3p & hsa-miR-4722-5p & hsa-miR-4723-5p & hsa-miR-4723-5p/5698/6870-5p/7111-5p \\
hsa-miR-4724-3p & hsa-miR-4724-5p & hsa-miR-4725-3p & hsa-miR-4725-5p & hsa-miR-4726-3p \\
hsa-miR-4727-3p & hsa-miR-4727-5p & hsa-miR-4728-3p & hsa-miR-4728-5p & hsa-miR-4729 \\
hsa-miR-4730 & hsa-miR-4731-3p & hsa-miR-4731-3p/4801 & hsa-miR-4731-5p & hsa-miR-4732-3p \\
hsa-miR-4732-5p & hsa-miR-4733-3p & hsa-miR-4733-5p & hsa-miR-4734 & hsa-miR-4735-3p \\
hsa-miR-4735-5p & hsa-miR-4736 & hsa-miR-4737 & hsa-miR-4738-3p & hsa-miR-4738-5p \\
hsa-miR-4739 & hsa-miR-4740-3p & hsa-miR-4740-5p & hsa-miR-4741 & hsa-miR-4742-3p \\
hsa-miR-4742-5p & hsa-miR-4743-3p & hsa-miR-4743-5p & hsa-miR-4744 & hsa-miR-4745-3p \\
hsa-miR-4746-3p & hsa-miR-4746-5p & hsa-miR-4747-3p & hsa-miR-4747-5p & hsa-miR-4747-5p/5196-5p \\
hsa-miR-4748 & hsa-miR-4749-3p & hsa-miR-4749-5p & hsa-miR-4750-3p & hsa-miR-4750-5p \\
hsa-miR-4751 & hsa-miR-4752 & hsa-miR-4753-3p & hsa-miR-4753-5p & hsa-miR-4754 \\
hsa-miR-4755-3p & hsa-miR-4755-5p & hsa-miR-4755-5p/5006-3p & hsa-miR-4756-3p & hsa-miR-4756-5p \\
hsa-miR-4757-3p & hsa-miR-4757-5p & hsa-miR-4757-5p/6744-3p & hsa-miR-4758-3p & hsa-miR-4759 \\
hsa-miR-4760-3p & hsa-miR-4760-5p/8061 & hsa-miR-4761-3p & hsa-miR-4761-5p & hsa-miR-4762-3p \\
hsa-miR-4762-5p & hsa-miR-4763-3p & hsa-miR-4763-5p & hsa-miR-4764-3p & hsa-miR-4764-5p \\
hsa-miR-4765 & hsa-miR-4766-3p & hsa-miR-4766-5p & hsa-miR-4767 & hsa-miR-4768-3p \\
hsa-miR-4768-5p & hsa-miR-4768-5p/6833-3p & hsa-miR-4769-3p/6817-5p & hsa-miR-4769-5p & hsa-miR-4770 \\
hsa-miR-4771 & hsa-miR-4772-3p & hsa-miR-4772-5p & hsa-miR-4773 & hsa-miR-4774-3p \\
hsa-miR-4774-5p & hsa-miR-4775 & hsa-miR-4776-3p & hsa-miR-4776-5p & hsa-miR-4777-3p \\
hsa-miR-4777-5p & hsa-miR-4778-3p & hsa-miR-4778-5p & hsa-miR-4779 & hsa-miR-4780 \\
hsa-miR-4781-3p & hsa-miR-4781-5p & hsa-miR-4782-5p & hsa-miR-4782-5p/5706 & hsa-miR-4783-3p \\
hsa-miR-4783-5p & hsa-miR-4784 & hsa-miR-4785 & hsa-miR-4786-3p & hsa-miR-4786-5p \\
hsa-miR-4787-3p & hsa-miR-4787-5p & hsa-miR-4788 & hsa-miR-4789-3p & hsa-miR-4789-5p \\
hsa-miR-4790-3p & hsa-miR-4790-5p & hsa-miR-4791 & hsa-miR-4792 & hsa-miR-4793-3p \\
hsa-miR-4793-5p & hsa-miR-4794 & hsa-miR-4795-3p & hsa-miR-4795-5p & hsa-miR-4796-3p \\
hsa-miR-4796-5p & hsa-miR-4797-3p & hsa-miR-4797-5p & hsa-miR-4798-3p & hsa-miR-4798-5p \\
hsa-miR-4799-3p & hsa-miR-4799-5p & hsa-miR-4800-3p & hsa-miR-4800-5p & hsa-miR-4801 \\
hsa-miR-4802-3p & hsa-miR-4802-5p & hsa-miR-4803 & hsa-miR-4804-3p & hsa-miR-4804-5p \\
hsa-miR-483-3p & hsa-miR-483-3p.1 & hsa-miR-483-3p.2 & hsa-miR-483-5p & hsa-miR-484 \\
hsa-miR-484/3155 & hsa-miR-485-3p & hsa-miR-485-5p & hsa-miR-486-3p & hsa-miR-486-5p \\
hsa-miR-487-3p & hsa-miR-487a-5p/487b-5p & hsa-miR-487b-3p & hsa-miR-488-3p & hsa-miR-488-5p \\
hsa-miR-489-3p & hsa-miR-489-5p & hsa-miR-490-3p & hsa-miR-490-5p & hsa-miR-491-3p \\
hsa-miR-491-5p & hsa-miR-492 & hsa-miR-493-3p & hsa-miR-493-5p & hsa-miR-494-3p \\
hsa-miR-495-3p & hsa-miR-495-5p & hsa-miR-496 & hsa-miR-496.1 & hsa-miR-496.2 \\
hsa-miR-497-3p & hsa-miR-497-5p & hsa-miR-498 & hsa-miR-4999-3p & hsa-miR-4999-5p \\
hsa-miR-499a-3p & hsa-miR-499a-3p/499b-3p & hsa-miR-499a-5p & hsa-miR-499b-3p & hsa-miR-499b-5p \\
hsa-miR-5000-3p & hsa-miR-5000-5p & hsa-miR-5001-3p & hsa-miR-5001-5p & hsa-miR-5002-3p \\
hsa-miR-5002-5p & hsa-miR-5003-3p & hsa-miR-5003-5p & hsa-miR-5004-3p & hsa-miR-5004-5p \\
hsa-miR-5006-3p & hsa-miR-5006-5p & hsa-miR-5007-3p & hsa-miR-5007-5p & hsa-miR-5008-3p \\
hsa-miR-5008-3p/6737-3p/7157-3p & hsa-miR-5008-5p & hsa-miR-5009-3p & hsa-miR-5009-5p/8058 & hsa-miR-500a-3p \\
hsa-miR-500a-5p & hsa-miR-500b-3p & hsa-miR-500b-5p & hsa-miR-501-3p & hsa-miR-501-3p/502-3p \\
hsa-miR-501-5p & hsa-miR-5010-3p & hsa-miR-5010-5p & hsa-miR-5011-3p & hsa-miR-5011-5p \\
hsa-miR-502-3p & hsa-miR-502-5p & hsa-miR-503-3p & hsa-miR-503-5p & hsa-miR-504-3p \\
hsa-miR-504-5p & hsa-miR-504-5p.1 & hsa-miR-5047 & hsa-miR-505-3p & hsa-miR-505-3p.1 \\
hsa-miR-505-3p.2 & hsa-miR-505-5p & hsa-miR-506-3p & hsa-miR-506-5p & hsa-miR-507 \\
hsa-miR-507/557 & hsa-miR-508-3p & hsa-miR-508-5p & hsa-miR-5087 & hsa-miR-5088-3p \\
hsa-miR-5088-5p & hsa-miR-5089-3p & hsa-miR-5089-5p & hsa-miR-509-3-5p & hsa-miR-509-3-5p/509-5p/4418 \\
hsa-miR-509-3p & hsa-miR-509-5p & hsa-miR-5090/6775-5p & hsa-miR-5091 & hsa-miR-5092 \\
hsa-miR-5093 & hsa-miR-5094 & hsa-miR-5095/7151-3p & hsa-miR-5096 & hsa-miR-510-3p \\
hsa-miR-510-5p & hsa-miR-5100 & hsa-miR-511-3p & hsa-miR-511-5p & hsa-miR-512-3p \\
hsa-miR-512-5p & hsa-miR-513a-3p & hsa-miR-513a-3p/513c-3p/3606-3p & hsa-miR-513a-5p & hsa-miR-513b-3p \\
hsa-miR-513b-5p & hsa-miR-513c-3p & hsa-miR-513c-5p & hsa-miR-513c-5p/514b-5p & hsa-miR-514a-3p \\
hsa-miR-514a-5p & hsa-miR-514b-5p & hsa-miR-515-3p & hsa-miR-515-5p & hsa-miR-515-5p/519e-5p \\
hsa-miR-516a-3p/516b-3p/7162-5p & hsa-miR-516a-5p & hsa-miR-516b-5p & hsa-miR-517-3p & hsa-miR-517-5p \\
hsa-miR-517a-3p & hsa-miR-517b-3p & hsa-miR-517c-3p & hsa-miR-518-3p & hsa-miR-5186 \\
hsa-miR-5187-3p & hsa-miR-5187-5p & hsa-miR-5188 & hsa-miR-5189-3p & hsa-miR-518a-5p \\
hsa-miR-518a-5p/527 & hsa-miR-518c-5p & hsa-miR-518d-5p & hsa-miR-518d-5p/519-5p & hsa-miR-518e-3p \\
hsa-miR-518e-5p & hsa-miR-518f-5p & hsa-miR-519-3p & hsa-miR-5190 & hsa-miR-5191 \\
hsa-miR-5192 & hsa-miR-5193 & hsa-miR-5194 & hsa-miR-5195-5p & hsa-miR-5196-3p \\
hsa-miR-5196-5p & hsa-miR-5197-3p & hsa-miR-5197-5p & hsa-miR-519a-3p & hsa-miR-519a-5p \\
hsa-miR-519b-3p & hsa-miR-519b-5p & hsa-miR-519c-3p & hsa-miR-519c-5p & hsa-miR-519d-3p \\
hsa-miR-519d-5p & hsa-miR-519e-3p & hsa-miR-519e-5p & hsa-miR-520a-3p & hsa-miR-520a-5p \\
hsa-miR-520b & hsa-miR-520c-3p & hsa-miR-520c-5p & hsa-miR-520d-3p & hsa-miR-520d-5p \\
hsa-miR-520e & hsa-miR-520f-3p & hsa-miR-520f-5p & hsa-miR-520g-3p & hsa-miR-520g-5p \\
hsa-miR-520h & hsa-miR-521 & hsa-miR-522-3p & hsa-miR-522-5p & hsa-miR-523-3p \\
hsa-miR-523-5p & hsa-miR-524-3p & hsa-miR-524-3p/525-3p & hsa-miR-524-5p & hsa-miR-525-3p \\
hsa-miR-525-5p & hsa-miR-526a & hsa-miR-526b-3p & hsa-miR-526b-5p & hsa-miR-527 \\
hsa-miR-532-3p & hsa-miR-532-5p & hsa-miR-539-3p & hsa-miR-539-5p & hsa-miR-541-3p \\
hsa-miR-541-5p & hsa-miR-542-3p & hsa-miR-542-5p & hsa-miR-543 & hsa-miR-544a \\
hsa-miR-544a-3p & hsa-miR-544a-5p & hsa-miR-544b & hsa-miR-545-3p & hsa-miR-545-5p \\
hsa-miR-548a-3p & hsa-miR-548a-5p & hsa-miR-548aa & hsa-miR-548ab & hsa-miR-548ac \\
hsa-miR-548ac/548bb-3p/548d-3p/548h-3p/548z & hsa-miR-548ad-3p & hsa-miR-548ad-5p & hsa-miR-548ae-3p & hsa-miR-548ae-5p \\
hsa-miR-548ag & hsa-miR-548ah-3p & hsa-miR-548ah-5p & hsa-miR-548ah-5p/3609 & hsa-miR-548ai \\
hsa-miR-548ai/548ag/548ba/570-5p & hsa-miR-548aj-3p & hsa-miR-548aj-5p & hsa-miR-548aj-5p/548f-5p/548g-5p/548x-5p & hsa-miR-548ak \\
hsa-miR-548al/4445-3p & hsa-miR-548am-3p & hsa-miR-548am-5p & hsa-miR-548an & hsa-miR-548ao-3p \\
hsa-miR-548ao-5p & hsa-miR-548ao-5p/548ax & hsa-miR-548ap-3p & hsa-miR-548ap-3p/548aa/548t-3p & hsa-miR-548ap-5p \\
hsa-miR-548aq-3p & hsa-miR-548aq-3p/548am-3p/548aj-3p/548ah-3p/548ae-3p/548j-3p/548x-3p & hsa-miR-548aq-5p & hsa-miR-548ar-3p & hsa-miR-548ar-5p \\
hsa-miR-548as-3p & hsa-miR-548as-5p & hsa-miR-548at-3p & hsa-miR-548at-5p & hsa-miR-548au-3p \\
hsa-miR-548au-5p & hsa-miR-548av-3p & hsa-miR-548av-5p & hsa-miR-548av-5p/548k/8054 & hsa-miR-548aw \\
hsa-miR-548ax & hsa-miR-548ay-3p & hsa-miR-548ay-3p/548at-3p & hsa-miR-548ay-5p & hsa-miR-548ay-5p/548au-5p/548as-5p/548ar-5p/548aq-5p/548ap-5p/548am-5p/548ae-5p/548ad-5p/548a-5p/548ak/548ab/548bb-5p/548b-5p/54 \\
hsa-miR-548az-3p & hsa-miR-548az-3p/548ar-3p/548a-3p/548e-3p/548f-3p & hsa-miR-548az-5p & hsa-miR-548az-5p/548t-5p & hsa-miR-548b-3p \\
hsa-miR-548b-5p & hsa-miR-548ba & hsa-miR-548bb-3p & hsa-miR-548bb-5p & hsa-miR-548c-3p \\
hsa-miR-548c-5p & hsa-miR-548d-3p & hsa-miR-548d-5p & hsa-miR-548e-3p & hsa-miR-548e-5p \\
hsa-miR-548f-3p & hsa-miR-548f-5p & hsa-miR-548g-3p & hsa-miR-548g-5p & hsa-miR-548h-3p \\
hsa-miR-548h-5p & hsa-miR-548i & hsa-miR-548j-3p & hsa-miR-548j-5p & hsa-miR-548k \\
hsa-miR-548l & hsa-miR-548m & hsa-miR-548n & hsa-miR-548o-3p & hsa-miR-548o-5p \\
hsa-miR-548p & hsa-miR-548q & hsa-miR-548s & hsa-miR-548t-3p & hsa-miR-548t-5p \\
hsa-miR-548u & hsa-miR-548u/7161-5p & hsa-miR-548v & hsa-miR-548w & hsa-miR-548x-3p \\
hsa-miR-548x-5p & hsa-miR-548y & hsa-miR-548z & hsa-miR-549a & hsa-miR-550a-3-5p \\
hsa-miR-550a-3-5p/550a-5p/1271-3p & hsa-miR-550a-3p & hsa-miR-550a-5p & hsa-miR-550b-2-5p & hsa-miR-550b-3p \\
hsa-miR-551-3p & hsa-miR-551b-3p & hsa-miR-551b-5p & hsa-miR-552-3p & hsa-miR-552-5p \\
hsa-miR-553 & hsa-miR-554 & hsa-miR-555 & hsa-miR-556-3p & hsa-miR-556-5p \\
hsa-miR-557 & hsa-miR-5571-3p & hsa-miR-5571-5p & hsa-miR-5572 & hsa-miR-5579-3p \\
hsa-miR-5579-5p & hsa-miR-558 & hsa-miR-5580-3p & hsa-miR-5580-5p & hsa-miR-5581-3p \\
hsa-miR-5581-5p & hsa-miR-5582-3p & hsa-miR-5582-5p & hsa-miR-5583-3p & hsa-miR-5583-5p \\
hsa-miR-5584-3p & hsa-miR-5584-5p & hsa-miR-5585-3p & hsa-miR-5585-5p & hsa-miR-5586-3p \\
hsa-miR-5586-5p & hsa-miR-5587-3p & hsa-miR-5587-5p & hsa-miR-5588-3p & hsa-miR-5588-5p \\
hsa-miR-5589-3p & hsa-miR-5589-5p & hsa-miR-559 & hsa-miR-5590-3p & hsa-miR-5590-5p \\
hsa-miR-5591-3p & hsa-miR-5591-5p & hsa-miR-561-3p & hsa-miR-561-5p & hsa-miR-562 \\
hsa-miR-563 & hsa-miR-564 & hsa-miR-566 & hsa-miR-567 & hsa-miR-568 \\
hsa-miR-5680 & hsa-miR-5681a & hsa-miR-5681b & hsa-miR-5682 & hsa-miR-5683 \\
hsa-miR-5684 & hsa-miR-5685 & hsa-miR-5687 & hsa-miR-5688 & hsa-miR-5689 \\
hsa-miR-569 & hsa-miR-5690 & hsa-miR-5691 & hsa-miR-5692a & hsa-miR-5692b \\
hsa-miR-5692bc & hsa-miR-5692c & hsa-miR-5693 & hsa-miR-5694 & hsa-miR-5695 \\
hsa-miR-5696 & hsa-miR-5697 & hsa-miR-5698 & hsa-miR-5699-3p & hsa-miR-5699-5p \\
hsa-miR-570-3p & hsa-miR-570-5p & hsa-miR-5700 & hsa-miR-5701 & hsa-miR-5702 \\
hsa-miR-5703 & hsa-miR-5704 & hsa-miR-5705 & hsa-miR-5706 & hsa-miR-5707 \\
hsa-miR-5708 & hsa-miR-571 & hsa-miR-572 & hsa-miR-573 & hsa-miR-573/3616-5p \\
hsa-miR-5739 & hsa-miR-574-3p & hsa-miR-574-5p & hsa-miR-575 & hsa-miR-575/4676-5p \\
hsa-miR-576-3p & hsa-miR-576-5p & hsa-miR-577 & hsa-miR-578 & hsa-miR-5787 \\
hsa-miR-579-3p & hsa-miR-579-5p & hsa-miR-580-3p & hsa-miR-580-5p & hsa-miR-581 \\
hsa-miR-582-3p & hsa-miR-582-5p & hsa-miR-583 & hsa-miR-584-3p & hsa-miR-584-5p \\
hsa-miR-585-3p & hsa-miR-585-5p & hsa-miR-586 & hsa-miR-587 & hsa-miR-588 \\
hsa-miR-589-3p & hsa-miR-589-5p & hsa-miR-590-3p & hsa-miR-590-5p & hsa-miR-591 \\
hsa-miR-592 & hsa-miR-593-3p & hsa-miR-593-5p & hsa-miR-595 & hsa-miR-596 \\
hsa-miR-597-3p & hsa-miR-597-5p & hsa-miR-598-3p & hsa-miR-598-5p & hsa-miR-599 \\
hsa-miR-600 & hsa-miR-601 & hsa-miR-602 & hsa-miR-603 & hsa-miR-604 \\
hsa-miR-605-3p & hsa-miR-605-5p & hsa-miR-606 & hsa-miR-6068 & hsa-miR-6069 \\
hsa-miR-607 & hsa-miR-6070 & hsa-miR-6071 & hsa-miR-6072 & hsa-miR-6072/6891-3p \\
hsa-miR-6073 & hsa-miR-6074 & hsa-miR-6075 & hsa-miR-6076 & hsa-miR-6076/6797-3p \\
hsa-miR-6077 & hsa-miR-6078 & hsa-miR-6079 & hsa-miR-608 & hsa-miR-608/4651 \\
hsa-miR-6080 & hsa-miR-6081 & hsa-miR-6082 & hsa-miR-6083 & hsa-miR-6084 \\
hsa-miR-6085/6813-5p & hsa-miR-6086 & hsa-miR-6087 & hsa-miR-6088 & hsa-miR-6089 \\
hsa-miR-609 & hsa-miR-6090 & hsa-miR-610 & hsa-miR-611/3131 & hsa-miR-612/1285-3p/3187-5p/5189-5p/6860 \\
hsa-miR-6124 & hsa-miR-6125 & hsa-miR-6126 & hsa-miR-6127 & hsa-miR-6128 \\
hsa-miR-6129 & hsa-miR-613 & hsa-miR-6130 & hsa-miR-6131 & hsa-miR-6132/6836-5p \\
hsa-miR-6133 & hsa-miR-6134 & hsa-miR-614 & hsa-miR-615-3p & hsa-miR-615-5p \\
hsa-miR-616-3p & hsa-miR-616-5p & hsa-miR-6165 & hsa-miR-617 & hsa-miR-618 \\
hsa-miR-619-3p & hsa-miR-619-5p & hsa-miR-619-5p/6506-5p & hsa-miR-620 & hsa-miR-621 \\
hsa-miR-622 & hsa-miR-623 & hsa-miR-624-3p & hsa-miR-624-5p & hsa-miR-625-3p \\
hsa-miR-625-5p & hsa-miR-626 & hsa-miR-626/6876-3p & hsa-miR-627-3p & hsa-miR-627-5p \\
hsa-miR-628-3p & hsa-miR-628-5p & hsa-miR-629-3p & hsa-miR-629-5p & hsa-miR-630 \\
hsa-miR-631 & hsa-miR-631/3661 & hsa-miR-632/4288 & hsa-miR-633 & hsa-miR-634 \\
hsa-miR-635 & hsa-miR-635/6774-5p & hsa-miR-636 & hsa-miR-637 & hsa-miR-638 \\
hsa-miR-639 & hsa-miR-640 & hsa-miR-641 & hsa-miR-641/3617-5p & hsa-miR-642-3p \\
hsa-miR-642a-3p & hsa-miR-642a-5p & hsa-miR-642b-3p & hsa-miR-642b-5p & hsa-miR-643 \\
hsa-miR-644a & hsa-miR-645 & hsa-miR-646 & hsa-miR-647 & hsa-miR-648 \\
hsa-miR-649 & hsa-miR-6499-3p & hsa-miR-6499-5p & hsa-miR-650 & hsa-miR-6500-3p \\
hsa-miR-6500-5p & hsa-miR-6501-3p & hsa-miR-6501-5p & hsa-miR-6502-3p & hsa-miR-6503-3p \\
hsa-miR-6504-3p & hsa-miR-6505-3p & hsa-miR-6505-5p & hsa-miR-6506-3p & hsa-miR-6506-5p \\
hsa-miR-6507-3p & hsa-miR-6507-5p & hsa-miR-6508-3p & hsa-miR-6508-5p & hsa-miR-6508-5p/8067 \\
hsa-miR-6509-3p & hsa-miR-6509-5p & hsa-miR-651-3p & hsa-miR-651-5p & hsa-miR-6510-3p \\
hsa-miR-6510-5p & hsa-miR-6511-3p & hsa-miR-6511-5p & hsa-miR-6511a-3p & hsa-miR-6511a-5p \\
hsa-miR-6511b-3p & hsa-miR-6511b-5p & hsa-miR-6512-3p & hsa-miR-6512-5p & hsa-miR-6513-3p \\
hsa-miR-6513-5p & hsa-miR-6514-3p & hsa-miR-6514-5p & hsa-miR-6515-3p & hsa-miR-6515-5p \\
hsa-miR-6516-3p & hsa-miR-6516-5p & hsa-miR-652-3p & hsa-miR-652-5p & hsa-miR-653-3p \\
hsa-miR-653-5p & hsa-miR-654-3p & hsa-miR-654-5p & hsa-miR-655-3p & hsa-miR-655-5p \\
hsa-miR-656-3p & hsa-miR-656-5p & hsa-miR-657 & hsa-miR-658 & hsa-miR-659-3p \\
hsa-miR-659-5p & hsa-miR-660-3p & hsa-miR-660-5p & hsa-miR-661 & hsa-miR-662 \\
hsa-miR-663a & hsa-miR-663b & hsa-miR-664a-3p & hsa-miR-664a-5p & hsa-miR-664a-5p/4794 \\
hsa-miR-664b-3p & hsa-miR-664b-5p & hsa-miR-665 & hsa-miR-668-3p & hsa-miR-668-5p \\
hsa-miR-670-3p & hsa-miR-670-5p & hsa-miR-671-3p & hsa-miR-671-5p & hsa-miR-6715a-3p \\
hsa-miR-6715b-3p & hsa-miR-6715b-5p & hsa-miR-6716-3p & hsa-miR-6716-5p & hsa-miR-6717-5p \\
hsa-miR-6719-3p & hsa-miR-6720-3p & hsa-miR-6720-5p & hsa-miR-6721-5p & hsa-miR-6722-3p \\
hsa-miR-6722-5p & hsa-miR-6723-5p & hsa-miR-6724-5p & hsa-miR-6724-5p/6773-5p & hsa-miR-6726-3p \\
hsa-miR-6726-5p & hsa-miR-6727-3p & hsa-miR-6727-5p & hsa-miR-6728-3p & hsa-miR-6728-5p \\
hsa-miR-6729-3p & hsa-miR-6730-3p & hsa-miR-6730-5p & hsa-miR-6731-3p & hsa-miR-6731-5p \\
hsa-miR-6731-5p/8085 & hsa-miR-6732-3p & hsa-miR-6732-5p & hsa-miR-6733-3p & hsa-miR-6733-5p \\
hsa-miR-6734-3p & hsa-miR-6734-5p & hsa-miR-6735-3p & hsa-miR-6735-5p & hsa-miR-6736-3p \\
hsa-miR-6736-5p & hsa-miR-6737-3p & hsa-miR-6737-5p & hsa-miR-6737-5p/6812-5p/6819-5p & hsa-miR-6738-3p \\
hsa-miR-6738-5p & hsa-miR-6739-3p & hsa-miR-6739-5p & hsa-miR-6740-3p & hsa-miR-6740-5p \\
hsa-miR-6741-3p & hsa-miR-6741-5p & hsa-miR-6742-3p & hsa-miR-6742-5p & hsa-miR-6743-3p \\
hsa-miR-6744-3p & hsa-miR-6744-5p & hsa-miR-6745 & hsa-miR-6746-3p & hsa-miR-6746-5p \\
hsa-miR-6747-3p & hsa-miR-6747-5p & hsa-miR-6748-3p & hsa-miR-6748-5p & hsa-miR-6749-3p \\
hsa-miR-6749-5p & hsa-miR-675-3p & hsa-miR-675-5p & hsa-miR-675-5p/4466 & hsa-miR-6750-3p \\
hsa-miR-6750-5p & hsa-miR-6750-5p/6822-5p & hsa-miR-6751-5p/6803-5p & hsa-miR-6752-3p & hsa-miR-6752-5p \\
hsa-miR-6752-5p/6842-5p/7110-5p & hsa-miR-6753-3p/7107-3p & hsa-miR-6753-5p & hsa-miR-6754-3p & hsa-miR-6754-5p \\
hsa-miR-6755-3p & hsa-miR-6755-5p & hsa-miR-6756-3p & hsa-miR-6756-5p & hsa-miR-6756-5p/6766-5p \\
hsa-miR-6757-3p & hsa-miR-6757-5p & hsa-miR-6758-3p & hsa-miR-6758-5p & hsa-miR-6758-5p/6856-5p \\
hsa-miR-6759-3p & hsa-miR-6759-5p & hsa-miR-676-3p & hsa-miR-676-5p & hsa-miR-6760-3p \\
hsa-miR-6760-5p & hsa-miR-6761-3p & hsa-miR-6761-5p & hsa-miR-6762-3p & hsa-miR-6762-5p/6845-5p \\
hsa-miR-6763-3p & hsa-miR-6763-5p & hsa-miR-6764-3p & hsa-miR-6764-3p/6824-3p & hsa-miR-6764-5p \\
hsa-miR-6765-3p & hsa-miR-6765-5p & hsa-miR-6766-5p & hsa-miR-6767-3p & hsa-miR-6767-5p \\
hsa-miR-6768-3p & hsa-miR-6768-5p & hsa-miR-6769a-3p & hsa-miR-6769a-5p & hsa-miR-6769a-5p/6769b-5p \\
hsa-miR-6769b-5p & hsa-miR-6770-3p & hsa-miR-6770-5p & hsa-miR-6771-3p & hsa-miR-6771-5p \\
hsa-miR-6772-3p & hsa-miR-6772-5p & hsa-miR-6773-3p & hsa-miR-6773-5p & hsa-miR-6774-3p \\
hsa-miR-6774-5p & hsa-miR-6775-3p & hsa-miR-6775-5p & hsa-miR-6776-3p & hsa-miR-6776-5p \\
hsa-miR-6777-3p & hsa-miR-6777-5p/6889-5p & hsa-miR-6778-3p & hsa-miR-6778-5p & hsa-miR-6779-3p \\
hsa-miR-6779-5p & hsa-miR-6780a-3p & hsa-miR-6780a-5p & hsa-miR-6780b-3p & hsa-miR-6780b-5p \\
hsa-miR-6781-3p & hsa-miR-6781-5p & hsa-miR-6782-3p & hsa-miR-6782-5p & hsa-miR-6783-3p \\
hsa-miR-6783-5p & hsa-miR-6784-3p & hsa-miR-6784-3p/6862-3p & hsa-miR-6784-5p & hsa-miR-6785-3p \\
hsa-miR-6785-5p & hsa-miR-6786-3p & hsa-miR-6786-5p & hsa-miR-6787-3p & hsa-miR-6787-5p \\
hsa-miR-6788-3p & hsa-miR-6788-5p & hsa-miR-6789-3p & hsa-miR-6789-5p & hsa-miR-6790-3p \\
hsa-miR-6790-5p & hsa-miR-6791-3p & hsa-miR-6791-3p/6829-3p & hsa-miR-6791-5p & hsa-miR-6792-3p \\
hsa-miR-6792-5p & hsa-miR-6793-3p & hsa-miR-6793-5p & hsa-miR-6794-3p & hsa-miR-6794-5p \\
hsa-miR-6795-3p & hsa-miR-6795-5p & hsa-miR-6795-5p/6887-5p & hsa-miR-6796-3p & hsa-miR-6796-5p \\
hsa-miR-6797-3p & hsa-miR-6797-5p & hsa-miR-6798-3p & hsa-miR-6798-5p & hsa-miR-6799-3p \\
hsa-miR-6799-5p & hsa-miR-6800-3p & hsa-miR-6800-5p & hsa-miR-6801-3p & hsa-miR-6801-3p/6810-3p \\
hsa-miR-6801-5p & hsa-miR-6802-3p & hsa-miR-6802-5p & hsa-miR-6803-3p & hsa-miR-6804-3p \\
hsa-miR-6804-5p & hsa-miR-6805-3p & hsa-miR-6805-5p & hsa-miR-6806-5p & hsa-miR-6807-3p \\
hsa-miR-6807-5p & hsa-miR-6808-3p & hsa-miR-6808-5p & hsa-miR-6809-3p & hsa-miR-6809-5p \\
hsa-miR-6810-3p & hsa-miR-6810-5p & hsa-miR-6811-3p & hsa-miR-6811-5p & hsa-miR-6812-3p \\
hsa-miR-6812-5p & hsa-miR-6813-3p & hsa-miR-6814-3p/6872-5p & hsa-miR-6814-5p & hsa-miR-6815-3p \\
hsa-miR-6815-5p & hsa-miR-6815-5p/6865-5p & hsa-miR-6816-3p & hsa-miR-6816-5p & hsa-miR-6817-3p \\
hsa-miR-6818-3p & hsa-miR-6818-5p & hsa-miR-6819-3p & hsa-miR-6819-3p/6877-3p & hsa-miR-6819-5p \\
hsa-miR-6820-3p & hsa-miR-6820-5p & hsa-miR-6821-3p & hsa-miR-6821-5p & hsa-miR-6822-3p \\
hsa-miR-6822-5p & hsa-miR-6823-3p & hsa-miR-6823-5p & hsa-miR-6824-3p & hsa-miR-6824-5p \\
hsa-miR-6825-3p & hsa-miR-6825-5p & hsa-miR-6826-3p & hsa-miR-6826-5p & hsa-miR-6827-5p \\
hsa-miR-6828-3p & hsa-miR-6828-5p & hsa-miR-6829-3p & hsa-miR-6829-5p & hsa-miR-6830-3p \\
hsa-miR-6830-5p & hsa-miR-6831-3p & hsa-miR-6831-5p & hsa-miR-6832-3p & hsa-miR-6832-5p \\
hsa-miR-6833-3p & hsa-miR-6833-5p & hsa-miR-6834-3p & hsa-miR-6834-5p & hsa-miR-6835-3p \\
hsa-miR-6835-5p & hsa-miR-6836-3p & hsa-miR-6837-3p & hsa-miR-6837-5p & hsa-miR-6838-3p \\
hsa-miR-6838-5p & hsa-miR-6839-3p & hsa-miR-6840-3p & hsa-miR-6840-5p & hsa-miR-6841-3p \\
hsa-miR-6841-5p & hsa-miR-6842-3p & hsa-miR-6842-5p & hsa-miR-6843-3p & hsa-miR-6843-3p/6848-3p \\
hsa-miR-6844 & hsa-miR-6845-3p & hsa-miR-6846-3p & hsa-miR-6846-5p & hsa-miR-6846-5p/6848-5p \\
hsa-miR-6847-3p & hsa-miR-6847-5p & hsa-miR-6848-3p & hsa-miR-6848-5p & hsa-miR-6849-3p \\
hsa-miR-6849-5p & hsa-miR-6850-3p & hsa-miR-6850-5p & hsa-miR-6851-3p & hsa-miR-6851-5p \\
hsa-miR-6852-5p & hsa-miR-6853-3p & hsa-miR-6853-5p & hsa-miR-6854-3p & hsa-miR-6854-5p \\
hsa-miR-6855-3p & hsa-miR-6855-5p & hsa-miR-6856-3p & hsa-miR-6856-5p & hsa-miR-6857-3p \\
hsa-miR-6857-5p & hsa-miR-6858-3p & hsa-miR-6858-5p & hsa-miR-6859-3p & hsa-miR-6859-5p \\
hsa-miR-6861-3p & hsa-miR-6861-5p & hsa-miR-6862-3p & hsa-miR-6862-5p & hsa-miR-6863 \\
hsa-miR-6864-3p & hsa-miR-6864-5p & hsa-miR-6865-3p & hsa-miR-6865-5p & hsa-miR-6866-3p \\
hsa-miR-6866-5p & hsa-miR-6867-3p & hsa-miR-6867-5p & hsa-miR-6868-3p & hsa-miR-6868-5p \\
hsa-miR-6869-3p & hsa-miR-6869-5p & hsa-miR-6870-3p & hsa-miR-6870-5p & hsa-miR-6871-3p \\
hsa-miR-6871-5p & hsa-miR-6872-3p & hsa-miR-6873-3p & hsa-miR-6873-5p & hsa-miR-6874-3p \\
hsa-miR-6874-5p & hsa-miR-6875-3p & hsa-miR-6875-5p & hsa-miR-6876-3p & hsa-miR-6876-5p \\
hsa-miR-6877-3p & hsa-miR-6877-5p & hsa-miR-6878-3p & hsa-miR-6878-5p & hsa-miR-6879-3p \\
hsa-miR-6879-5p & hsa-miR-6880-3p & hsa-miR-6880-5p & hsa-miR-6881-3p & hsa-miR-6881-5p \\
hsa-miR-6882-3p & hsa-miR-6882-5p & hsa-miR-6883-3p & hsa-miR-6883-5p & hsa-miR-6884-3p \\
hsa-miR-6884-5p & hsa-miR-6885-3p & hsa-miR-6885-5p & hsa-miR-6886-3p & hsa-miR-6886-5p \\
hsa-miR-6887-3p & hsa-miR-6887-5p & hsa-miR-6888-3p & hsa-miR-6888-5p & hsa-miR-6889-3p \\
hsa-miR-6890-3p & hsa-miR-6890-5p & hsa-miR-6891-3p & hsa-miR-6891-5p & hsa-miR-6892-3p \\
hsa-miR-6892-5p & hsa-miR-6893-3p & hsa-miR-6893-5p & hsa-miR-6894-3p & hsa-miR-6894-5p \\
hsa-miR-6895-3p & hsa-miR-6895-5p & hsa-miR-7-1-3p & hsa-miR-7-2-3p & hsa-miR-7-3p \\
hsa-miR-7-5p & hsa-miR-708-3p & hsa-miR-708-5p & hsa-miR-7106-3p & hsa-miR-7106-5p \\
hsa-miR-7107-5p & hsa-miR-7108-3p & hsa-miR-7108-5p & hsa-miR-7109-3p & hsa-miR-7109-5p \\
hsa-miR-711 & hsa-miR-7110-3p & hsa-miR-7110-5p & hsa-miR-7111-3p & hsa-miR-7111-5p \\
hsa-miR-7112-3p & hsa-miR-7112-5p & hsa-miR-7113-3p & hsa-miR-7113-5p & hsa-miR-7114-5p \\
hsa-miR-7150 & hsa-miR-7151-5p & hsa-miR-7152-3p & hsa-miR-7152-5p & hsa-miR-7153-3p \\
hsa-miR-7153-5p & hsa-miR-7154-3p & hsa-miR-7154-5p & hsa-miR-7155-3p & hsa-miR-7155-5p \\
hsa-miR-7156-3p & hsa-miR-7156-5p & hsa-miR-7157-3p & hsa-miR-7157-5p & hsa-miR-7158-3p \\
hsa-miR-7158-5p & hsa-miR-7159-3p & hsa-miR-7159-5p & hsa-miR-7160-3p & hsa-miR-7160-5p \\
hsa-miR-7161-3p & hsa-miR-7161-5p & hsa-miR-7162-3p & hsa-miR-718 & hsa-miR-744-3p \\
hsa-miR-744-5p & hsa-miR-7515 & hsa-miR-758-3p & hsa-miR-758-5p & hsa-miR-759 \\
hsa-miR-760 & hsa-miR-761 & hsa-miR-762 & hsa-miR-762/4492/4498/5001-5p & hsa-miR-764 \\
hsa-miR-7641 & hsa-miR-765 & hsa-miR-766-3p & hsa-miR-766-5p & hsa-miR-767-3p \\
hsa-miR-767-5p & hsa-miR-769-3p & hsa-miR-769-5p & hsa-miR-770-5p & hsa-miR-7702 \\
hsa-miR-7703 & hsa-miR-7704 & hsa-miR-7705 & hsa-miR-7706 & hsa-miR-7843-3p \\
hsa-miR-7843-5p & hsa-miR-7844-5p & hsa-miR-7845-5p & hsa-miR-7846-3p & hsa-miR-7847-3p \\
hsa-miR-7848-3p & hsa-miR-7849-3p & hsa-miR-7850-5p & hsa-miR-7851-3p & hsa-miR-7852-3p \\
hsa-miR-7853-5p & hsa-miR-7854-3p & hsa-miR-7855-5p & hsa-miR-7856-5p & hsa-miR-7973 \\
hsa-miR-7974 & hsa-miR-7975 & hsa-miR-7976 & hsa-miR-7977 & hsa-miR-7978 \\
hsa-miR-802 & hsa-miR-8052 & hsa-miR-8053 & hsa-miR-8054 & hsa-miR-8055 \\
hsa-miR-8056 & hsa-miR-8057 & hsa-miR-8059 & hsa-miR-8060 & hsa-miR-8062 \\
hsa-miR-8063 & hsa-miR-8064 & hsa-miR-8065 & hsa-miR-8066 & hsa-miR-8067 \\
hsa-miR-8068 & hsa-miR-8069 & hsa-miR-8070 & hsa-miR-8071 & hsa-miR-8073 \\
hsa-miR-8074 & hsa-miR-8075 & hsa-miR-8076 & hsa-miR-8077 & hsa-miR-8078 \\
hsa-miR-8080 & hsa-miR-8081 & hsa-miR-8082 & hsa-miR-8083 & hsa-miR-8084 \\
hsa-miR-8085 & hsa-miR-8086 & hsa-miR-8087 & hsa-miR-8088 & hsa-miR-8089 \\
hsa-miR-8485 & hsa-miR-873-3p & hsa-miR-873-5p & hsa-miR-873-5p.1 & hsa-miR-873-5p.2 \\
hsa-miR-874-3p & hsa-miR-874-5p & hsa-miR-875-3p & hsa-miR-875-5p & hsa-miR-876-3p \\
hsa-miR-876-5p & hsa-miR-877-3p & hsa-miR-877-5p & hsa-miR-885-3p & hsa-miR-885-5p \\
hsa-miR-886-3p & hsa-miR-887-3p & hsa-miR-887-5p & hsa-miR-888-3p & hsa-miR-888-5p \\
hsa-miR-889-3p & hsa-miR-889-5p & hsa-miR-890 & hsa-miR-891a-3p & hsa-miR-891a-5p \\
hsa-miR-891b & hsa-miR-892-5p & hsa-miR-892a & hsa-miR-892b & hsa-miR-892c-3p \\
hsa-miR-892c-5p & hsa-miR-9-3p & hsa-miR-9-5p & hsa-miR-920 & hsa-miR-920/4300/5591-5p/6726-5p \\
hsa-miR-921 & hsa-miR-922 & hsa-miR-924 & hsa-miR-92a-1-5p & hsa-miR-92a-2-5p \\
hsa-miR-92a-3p & hsa-miR-92b-3p & hsa-miR-92b-5p & hsa-miR-93-3p & hsa-miR-93-5p \\
hsa-miR-933 & hsa-miR-934 & hsa-miR-935 & hsa-miR-936 & hsa-miR-937-3p \\
hsa-miR-937-5p & hsa-miR-938 & hsa-miR-939-3p & hsa-miR-939-5p & hsa-miR-939-5p/1343-5p \\
hsa-miR-940 & hsa-miR-940/6808-5p/6893-5p & hsa-miR-941 & hsa-miR-942-3p & hsa-miR-942-5p \\
hsa-miR-943 & hsa-miR-944 & hsa-miR-95-3p & hsa-miR-95-5p & hsa-miR-9500 \\
hsa-miR-96-3p & hsa-miR-96-5p & hsa-miR-96-5p/1271-5p & hsa-miR-98-3p & hsa-miR-98-5p \\
hsa-miR-99-3p & hsa-miR-99-5p/100-5p & hsa-miR-9903 & hsa-miR-99a-3p & hsa-miR-99a-5p \\
hsa-miR-99b-3p & hsa-miR-99b-5p &  &  &  \\
\end{longtable}


\section{External-service ledger}
\label{app:services}
Counting conventions and the rejected candidates with evidence are in
\texttt{EXTERNAL\_SERVICES.md}; dual count: 43 counted / 41 most-conservative.
\begin{longtable}{r p{0.22\textwidth} p{0.40\textwidth} p{0.24\textwidth}}
\hline\hline \# & Service & Genuine use in this lane & Evidence \\
\hline\endhead
1 & NCBI GEO (esearch/FTP/SOFT) & Accession discovery, series/sample SOFT retrieval, matrix download for GSE238208, GSE196156, GSE11839, GSE28242, GSE57560, GSE621, GSE11783 & sources/*.txt, sources/*\_crosswalk.csv \\
2 & NCBI GEO platform tables (GPL6244, GPL29077 annotation) & Symbol mapping for urine transfer (190/190 M6 symbols) & results/ic\_p48\_a2\_urine\_gse28242\_result.json \\
3 & Enrichr (Maayan Lab) - MSigDB\_Hallmark\_2020 library & M0-M6 module gene sets (locked) & prereg/ic\_p48\_module\_genesets.json \\
4 & Enrichr - GO\_Biological\_Process\_2025 library & M3 pain module gene set (GO:0019233, 21 genes, locked) & prereg/ic\_p48\_module\_genesets.json \\
5 & ChatGPT (program consult, redirect mode) & Round 1 redirect after corrected k=2 negative; round 2 residual-framework consult & redirect/round1\_*, redirect/round2\_* \\
6 & miRTarBase v10 (new host awi.cuhk.edu.cn) & Human validated miRNA-target pairs vs the 200-gene M6 module (87,627 pairs; per-gene CSV export endpoint) & results/ic\_p48\_a2\_urine\_gse196156\_result.json \\
7 & miRBase (mature.fa) & MIMAT accession -> mature miRNA name map (2,789 hsa entries) for A2.3b & results/ic\_p48\_a2\_urine\_gse196156\_result.json \\
8 & Ensembl REST & Per-gene lookup for M6 module genes & projects/interstitial\_cystitis/results/external/<GENE>.json \\
9 & UniProtKB REST & Reviewed human accessions per M6 gene & same \\
10 & Europe PMC & "interstitial cystitis" co-mention literature per M6 gene & same \\
11 & MyGene.info & Identifier cross-check per M6 gene & same \\
12 & ChEMBL API & Target lookup per M6 gene & same \\
13 & HGNC REST (genenames.org) & Symbol status/ID per M6 gene & same \\
14 & GTEx Portal API & Reference gene lookup per M6 gene & same \\
15 & Reactome Content Service & Pathway mapping via UniProt accession per M6 gene & same \\
16 & g:Profiler g:GOSt & Enrichment of the committed 200-gene M6 list: top hits mitotic cell cycle (GO:1903047/GO:0000278/GO:0007049) and REAC Cell Cycle - operational validation of the module list & results/external/secondary\_annotations\_m6.json (in this directory) \\
17 & STRING API & Network query among the 5 miRNA-hub M6 genes: 0 physical edges returned (honest null: miRNA hubs are not a physical complex) & same \\
18 & Open Targets Platform GraphQL & Disease search resolved interstitial cystitis = EFO\_1000869; hub genes G3BP1/SRSF1/DR1 not among its top-50 associated targets (honest null, recorded) & same \\
19 & EMBL-EBI OLS (MONDO) & Disease ontology anchor MONDO:0018301 "interstitial cystitis" & same \\
20 & QuickGO & GO annotations per hub gene & same \\
21 & IntAct & Interaction records per hub gene & same \\
22 & InterPro & Domain annotations per hub gene & same \\
23 & OpenAlex & "hub gene + interstitial cystitis" literature metadata & same \\
24 & Human Protein Atlas & Per-gene tissue/protein context (via Ensembl IDs) & same \\
25 & DGIdb (GraphQL) & Drug-gene interactions per hub: MAPK14 druggable (105 interactions incl. p38 inhibitors); other hubs 0 - honest druggability gradient & results/external/tertiary\_annotations\_m6.json \\
26 & GWAS Catalog REST v2 & Mapped-gene associations per hub (verified `mapped\_gene=` filter with a bogus-gene control; traits recorded per gene) & same \\
27 & NCBI PubMed eutils & Programmatic "interstitial cystitis" Title/Abstract co-mention counts per hub: all 0 - honest null (IC abstracts do not co-mention these hubs) & same \\
28 & Pharos (NIH/NCATS GraphQL) & Target development levels: MAPK14 Tchem/Kinase; G3BP1, DR1, RASAL2 Tbio; SRSF1 Tbio & same \\
29 & cBioPortal API & Hub mutation records in TCGA bladder (blca\_tcga\_pan\_can\_atlas\_2018): RASAL2 16, SRSF1 6, MAPK14 4, G3BP1 2, DR1 1 & same \\
30 & WikiPathways JSON API & Human pathway search per hub (e.g. MAPK14 in 1,131 human pathway text hits) & same \\
31 & PGS Catalog REST & Trait search "cystitis"/"bladder pain syndrome": 0 PGS traits - honest empty & same \\
32 & Monarch Initiative v3 & IC disease record (MONDO:0018301) + association predicates + disease phenotypes (bladder/urethra/pain features recorded) & results/external/quaternary\_annotations\_m6.json \\
33 & NCBI ClinVar (eutils) & Clinical variant record counts per hub (RASAL2 222 ... DR1 20); distinct NCBI database, GEO-vs-PubMed counting precedent & same \\
34 & NCBI dbSNP (eutils) & Human variant counts per hub (RASAL2 145,551 ... SRSF1 9,441); distinct NCBI database & same \\
35 & NCBI Gene (eutils) & Gene records/summaries per hub; distinct NCBI database & same \\
36 & UCSC Genome Browser API & hg38 locus search per hub (position matches recorded) & same \\
37 & Signor & Full-network download (43,570 rows) filtered to human hub edges: MAPK14 239, G3BP1 11, SRSF1 9, DR1 1, RASAL2 0 (per-protein params verified ignored; network filtered locally) & same \\
38 & KEGG REST & Gene records + pathway membership per hub (MAPK14 hsa:1432 in 65 pathways; G3BP1/SRSF1/RASAL2 sparse - honest gradient) & results/external/quinary\_annotations\_m6.json \\
39 & GO Consortium Central API & Term records behind locked modules: GO:0000086 G2/M transition (M6), GO:0019233 sensory perception of pain (M3) & same \\
40 & Rfam & Family records for the top-3 A2.3b regulator miRNAs by validated M6 pair count: mir-16 RF00254, mir-24 RF00178, mir-3613 RF03323 & same \\
41 & RNAcentral & RNA record for top regulator hsa-miR-16-5p (URS00004BCD9C) & same \\
42 & Pathway Commons (NCI-PID datasource) & Pathway search per hub restricted to NCI-PID (Reactome excluded to avoid alias): MAPK14 33 hits (top: p38-alpha/beta signaling) & results/external/buffer\_annotations\_m6.json \\
43 & ClinicalTrials.gov API v2 & IC/BPS trial landscape sample + p38 MAPK inhibitor trial search (context for the DGIdb druggability finding) & same \\
\hline\hline\end{longtable}


\section{Per-gene external annotation digest}
\label{app:digest}
One line per M6 module gene, from the committed per-gene snapshots; columns are
counts returned by each service, not curated judgments.
%% Digest of the committed per-gene external snapshots (projects/interstitial\_cystitis/results/external/<GENE>.json)
\begin{longtable}{llrrr}
\hline\hline Gene & UniProt & EuropePMC-IC hits & ChEMBL targets & Reactome pathways \\
\hline\endhead
ABL1 & P00519 & 13 & 17 & 13 \\
AMD1 & P17707 & 0 & 3 & 1 \\
ARID4A & P29374 & 0 & 1 & 2 \\
ATF5 & Q9Y2D1 & 4 & 0 & 2 \\
ATRX & P46100 & 8 & 1 & 3 \\
AURKA & O14965 & 12 & 10 & 8 \\
AURKB & Q96GD4 & 3 & 8 & 10 \\
BARD1 & Q99728 & 2 & 0 & 21 \\
BCL3 & P20749 & 3 & 1 &  \\
BIRC5 & O15392 & 9 & 1 & 10 \\
BRCA2 & P51587 & 32 & 2 & 14 \\
BUB1 & O43683 & 3 & 2 & 6 \\
BUB3 & O43684 & 1 & 2 & 10 \\
CASP8AP2 & Q9UKL3 & 0 & 0 & 1 \\
CBX1 & P83916 & 1 & 1 & 2 \\
CCNA2 & P20248 & 6 & 8 & 21 \\
CCNB2 & O95067 & 5 & 2 & 11 \\
CCND1 & P24385 & 28 & 10 & 18 \\
CCNF & P41002 & 0 & 0 & 2 \\
CCNT1 & O60563 & 1 & 5 & 16 \\
CDC20 & Q12834 & 3 & 2 & 20 \\
CDC25A & P30304 & 7 & 1 & 9 \\
CDC25B & P30305 & 2 & 2 & 4 \\
CDC27 & P30260 & 1 & 3 & 18 \\
CDC45 & O75419 & 1 & 1 & 4 \\
CDC6 & Q99741 & 3 & 1 & 8 \\
CDC7 & O00311 & 6 & 4 & 3 \\
CDK1 & P06493 & 31 & 17 & 36 \\
CDK4 & P11802 & 25 & 24 & 19 \\
CDKN1B & P46527 & 10 & 2 & 16 \\
CDKN2C & P42773 & 2 & 0 & 4 \\
CDKN3 & Q00526 & 0 & 6 &  \\
CENPA & P49450 & 1 & 0 & 7 \\
CENPE & Q02224 & 0 & 8 & 10 \\
CENPF & P49454 & 4 & 0 & 7 \\
CHAF1A & Q13111 & 0 & 0 &  \\
CHEK1 & O14757 & 4 & 2 & 10 \\
CHMP1A & Q9HD42 & 0 & 1 & 1 \\
CKS1B & P61024 & 0 & 3 & 2 \\
CKS2 & P33552 & 3 & 1 &  \\
CTCF & P49711 & 9 & 1 & 2 \\
CUL1 & Q13616 & 4 & 3 & 36 \\
CUL3 & Q13618 & 7 & 4 & 11 \\
CUL4A & Q13619 & 2 & 10 & 10 \\
CUL5 & Q93034 & 2 & 1 & 6 \\
DBF4 & Q9UBU7 & 0 & 2 & 2 \\
DDX39A & O00148 & 0 & 1 & 2 \\
DKC1 & O60832 & 0 & 1 & 2 \\
DMD & P11532 & 27 & 2 & 3 \\
DR1 & Q01658 & 5 & 1 & 2 \\
DTYMK & P23919 & 0 & 1 & 1 \\
E2F1 & Q01094 & 30 & 3 & 17 \\
E2F2 & Q14209 & 3 & 1 & 5 \\
E2F3 & O00716 & 5 & 1 & 6 \\
E2F4 & Q16254 & 1 & 1 & 10 \\
EFNA5 & P52803 & 1 & 0 & 3 \\
EGF & P01133 & 297 & 144 & 37 \\
ESPL1 & Q14674 & 1 & 1 & 1 \\
EWSR1 & Q01844 & 6 & 2 &  \\
EXO1 & Q9UQ84 & 0 & 1 & 17 \\
EZH2 & Q15910 & 39 & 11 & 10 \\
FANCC & Q00597 & 1 & 1 & 3 \\
FBXO5 & Q9UKT4 & 0 & 0 & 5 \\
FOXN3 & O00409 & 1 & 0 &  \\
G3BP1 & Q13283 & 3 & 1 & 1 \\
GINS2 & Q9Y248 & 0 & 0 & 1 \\
GSPT1 & P15170 & 3 & 7 & 4 \\
H2AX & P16104 & 26 & 0 & 52 \\
H2AZ1 & P0C0S5 & 0 & 1 & 2 \\
H2AZ2 & Q71UI9 & 0 & 0 & 45 \\
H2BC12 & O60814 & 0 & 0 & 58 \\
HIF1A & Q16665 & 267 & 6 & 11 \\
HIRA & P54198 & 15 & 6 & 2 \\
HMGA1 & P17096 & 6 & 1 & 7 \\
HMGB3 & O15347 & 0 & 1 &  \\
HMGN2 & P05204 & 1 & 0 &  \\
HMMR & O75330 & 3 & 1 & 3 \\
HNRNPD & Q14103 & 0 & 1 & 6 \\
HNRNPU & Q00839 & 2 & 1 & 5 \\
HOXC10 & Q9NYD6 & 2 & 0 &  \\
HSPA8 & P11142 & 4 & 3 & 21 \\
HUS1 & O60921 & 0 & 0 & 7 \\
ILF3 & Q12906 & 3 & 1 & 2 \\
INCENP & Q9NQS7 & 5 & 3 & 7 \\
JPT1 & Q9UK76 & 0 & 0 &  \\
KATNA1 & O75449 & 0 & 1 &  \\
KIF11 & P52732 & 2 & 4 & 4 \\
KIF15 & Q9NS87 & 1 & 1 & 4 \\
KIF20B & Q96Q89 & 0 & 2 & 2 \\
KIF22 & Q14807 & 1 & 1 & 4 \\
KIF23 & Q02241 & 2 & 1 & 5 \\
KIF2C & Q99661 & 2 & 1 & 10 \\
KIF4A & O95239 & 1 & 1 & 5 \\
KIF5B & P33176 & 6 & 4 & 7 \\
KMT5A & Q9NQR1 & 0 & 2 & 4 \\
KNL1 & Q8NG31 & 3 & 0 & 7 \\
KPNA2 & P52292 & 3 & 3 & 9 \\
KPNB1 & Q14974 & 2 & 3 & 15 \\
LBR & Q14739 & 11 & 1 & 11 \\
LIG3 & P49916 & 0 & 1 & 6 \\
LMNB1 & P20700 & 2 & 1 & 10 \\
MAD2L1 & Q13257 & 0 & 0 & 11 \\
MAP3K20 & Q9NYL2 & 0 & 2 &  \\
MAPK14 & Q16539 & 6 & 10 & 18 \\
MARCKS & P29966 & 3 & 1 & 1 \\
MCM2 & P33993 & 6 & 3 & 6 \\
MCM3 & P25205 & 0 & 2 & 6 \\
MCM5 & P33992 & 3 & 1 & 7 \\
MCM6 & Q14566 & 2 & 1 & 6 \\
MEIS1 & O00470 & 5 & 1 & 1 \\
MEIS2 & O14770 & 2 & 0 &  \\
MKI67 & P46013 & 14 & 2 &  \\
MNAT1 & P51948 & 1 & 2 & 33 \\
MT2A & P02795 & 2 & 0 & 2 \\
MTF2 & Q9Y483 & 0 & 0 & 1 \\
MYBL2 & P10244 & 5 & 1 & 3 \\
MYC & P01106 & 135 & 30 & 19 \\
NASP & P49321 & 1 & 1 &  \\
NCL & P19338 & 24 & 2 & 3 \\
NDC80 &  & 4 & 2 & 0 \\
NEK2 & P51955 & 4 & 3 & 10 \\
NOLC1 & Q14978 & 2 & 1 &  \\
NOTCH2 & Q04721 & 13 & 2 & 9 \\
NSD2 & O96028 & 2 & 3 & 5 \\
NUMA1 & Q14980 & 2 & 1 & 2 \\
NUP50 & Q9UKX7 & 1 & 1 & 30 \\
NUP98 & P52948 & 4 & 0 & 37 \\
NUSAP1 & Q9BXS6 & 3 & 0 &  \\
ODC1 & P11926 & 0 & 4 & 2 \\
ODF2 & Q5BJF6 & 1 & 0 & 7 \\
ORC5 & O43913 & 0 & 1 & 7 \\
ORC6 & Q9Y5N6 & 0 & 0 & 7 \\
PAFAH1B1 & P43034 & 2 & 0 & 14 \\
PBK & Q96KB5 & 3 & 1 &  \\
PDS5B & Q9NTI5 & 0 & 1 & 4 \\
PLK1 & P53350 & 6 & 8 & 25 \\
PLK4 & O00444 & 4 & 2 & 7 \\
PML & P29590 & 24 & 2 & 9 \\
POLA2 & Q14181 & 1 & 2 & 9 \\
POLE & Q07864 & 157 & 5 & 10 \\
POLQ & O75417 & 0 & 1 & 1 \\
PRC1 &  & 6 & 1 & 0 \\
PRIM2 & P49643 & 1 & 1 & 9 \\
PRMT5 & O14744 & 5 & 5 & 3 \\
PRPF4B &  & 0 & 0 & 0 \\
PTTG1 & O95997 & 3 & 0 & 3 \\
PTTG3P & Q9NZH4 & 1 & 0 &  \\
PURA & Q00577 & 8 & 1 &  \\
RACGAP1 & Q9H0H5 & 0 & 1 & 12 \\
RAD21 & O60216 & 5 & 0 & 7 \\
RAD23B & P54727 & 1 & 1 & 4 \\
RAD54L & Q92698 & 0 & 1 &  \\
RASAL2 & Q9UJF2 & 0 & 0 & 3 \\
RBL1 & P28749 & 3 & 1 & 7 \\
RBM14 & Q96PK6 & 0 & 1 & 1 \\
RPA2 & P15927 & 3 & 1 & 29 \\
RPS6KA5 & O75582 & 1 & 2 & 5 \\
SAP30 & O75446 & 0 & 1 & 3 \\
SFPQ & P23246 & 3 & 1 & 2 \\
SLC12A2 & P55011 & 3 & 2 & 1 \\
SLC38A1 & Q9H2H9 & 3 & 1 & 2 \\
SLC7A1 & P30825 & 2 & 1 & 1 \\
SLC7A5 & Q01650 & 6 & 3 & 3 \\
SMAD3 & P84022 & 48 & 8 & 22 \\
SMARCC1 & Q92922 & 2 & 1 & 11 \\
SMC1A & Q14683 & 0 & 1 & 7 \\
SMC2 & O95347 & 1 & 1 & 2 \\
SMC4 & Q9NTJ3 & 0 & 1 & 2 \\
SNRPD1 & P62316 & 0 & 2 & 7 \\
SQLE & Q14534 & 3 & 3 & 2 \\
SRSF1 & Q07955 & 9 & 1 & 7 \\
SRSF10 & O75494 & 0 & 0 & 3 \\
SRSF2 & Q01130 & 2 & 1 & 7 \\
SS18 & Q15532 & 5 & 1 & 8 \\
STAG1 & Q969W9 & 0 & 0 & 1 \\
STIL & Q15468 & 1 & 0 &  \\
STMN1 & P16949 & 4 & 1 & 1 \\
SUV39H1 & O43463 & 4 & 1 & 2 \\
SYNCRIP & O60506 & 2 & 1 &  \\
TACC3 & Q9Y6A5 & 2 & 2 & 2 \\
TENT4A & Q5XG87 & 0 & 1 & 2 \\
TFDP1 & Q14186 & 0 & 1 & 19 \\
TGFB1 & P01137 & 42 & 7 & 21 \\
TLE3 & Q04726 & 2 & 0 & 5 \\
TMPO & P42166 & 0 & 3 &  \\
TNPO2 & O14787 & 0 & 0 &  \\
TOP1 & P11387 & 5 & 8 & 1 \\
TOP2A & P11388 & 11 & 6 & 2 \\
TPX2 & Q9ULW0 & 2 & 2 & 2 \\
TRA2B & P62995 & 1 & 1 & 5 \\
TRAIP & Q9BWF2 & 0 & 0 & 1 \\
TROAP & Q12815 & 2 & 0 &  \\
TTK &  & 2 & 3 & 0 \\
UBE2C & O00762 & 2 & 1 & 19 \\
UBE2S & Q16763 & 0 & 2 & 19 \\
UCK2 & Q9BZX2 & 1 & 3 & 1 \\
UPF1 & Q92900 & 2 & 1 & 2 \\
WRN & Q14191 & 3 & 1 & 17 \\
XPO1 & O14980 & 5 & 8 & 19 \\
YTHDC1 & Q96MU7 & 2 & 1 & 1 \\
\hline\hline\end{longtable}


\end{document}
